\documentclass{article}
\usepackage[margin=1in]{geometry}

\setcounter{secnumdepth}{3}
% Useful packages
\usepackage{graphicx}
\usepackage{amsmath}
\usepackage{amssymb}
\usepackage{hyperref}
\usepackage{braket}
\usepackage{dsfont}
%\usepackage{showlabels}
\usepackage{amsthm}
\newtheorem{theorem}{Theorem}
\newtheorem{lemma}[theorem]{Lemma}
\newtheorem{corollary}[theorem]{Corollary}
\newtheorem{proposition}[theorem]{Proposition}
\theoremstyle{definition}
\newtheorem{definition}{Definition}[section]
\newtheorem{prop}{Conjecture}
\newtheorem{cor}{Corollary}
%-------------------packages-------------------------
\usepackage{enumerate}
\usepackage[utf8]{inputenc}
\usepackage{amsfonts}
\usepackage{amsbsy, mathtools}
% \usepackage{braket} % already loaded above
% \usepackage{graphicx} % already loaded above
\usepackage{xcolor}
%\usepackage{multirow}
% \usepackage{hyperref} % already loaded above
%\usepackage{caption}
%\usepackage{subcaption}
%\captionsetup{compatibility=false}
\usepackage{natbib}
%\usepackage{dcolumn}% Align table columns on decimal point
\usepackage{bm}% bold math
\usepackage{soul}
\usepackage{showlabels}

%Includes "References" in the table of contents
\usepackage[nottoc]{tocbibind}

\usepackage{bbm}
%\newtheorem{theorem}{Theorem}
%\newtheorem{lemma}[theorem]{Lemma}
%\newtheorem{corollary}[theorem]{Corollary}
%\theoremstyle{definition}
%\newtheorem{definition}{Definition}[section]
\usepackage{tikz}
\usepackage{quantikz}
\usepackage{textcomp}
\usepackage{caption}
\usepackage{microtype}
\newcommand{\Tr}{\text{tr}}
\newcommand{\mA}{\mathcal{A}}
\newcommand{\mB}{\mathcal{B}}
\newcommand{\nn}{\nonumber}
\newcommand{\mL}{\mathcal{L}}
\newcommand{\lb}{\left(}
\newcommand{\rb}{\right)}
\newcommand{\mU}{\mathcal{U}}
\newcommand{\mH}{\mathcal{H}}
\usepackage[T1]{fontenc}
\usepackage{lmodern}
\usepackage{subcaption}
\newcommand{\tOO}{\tilde{\mathcal{O}}}
\newcommand{\mO}{\mathcal{O}}
\newcommand{\mK}{\mathcal{K} }
\newcommand{\ep}{\epsilon}
\newcommand{\mbR}{\mathbb{R}}
\newcommand{\mM}{\mathcal{M}}
\newcommand{\mT}{\mathcal{T}}
\newcommand{\NL}[1]{\textcolor{blue}{#1}}
\newcommand{\p}{\partial}
\newcommand{\tr}{\text{tr}}
\newcommand{\tT}{\tilde{T}}

\usepackage{graphicx} % Required for inserting images

\title{Decay of correlations in ensemble-averaging theories}
\author{Nima Lashkari}
\date{January 2026}

\begin{document}
    \maketitle
    \tableofcontents

    \section*{Results}

\begin{enumerate}
    \item  {\bf Time-independent infinite-temperature ensemble:} We find that in time-independent ensemble-averaged quantum theories with a large
    $N$ limit and at infinite temperature, the $k$-twirl map becomes
    \begin{eqnarray}
        \mathcal{E}_g\lb U^{\otimes k}\otimes (U^*)^{\otimes k} \rb=e^{-t^2\mathcal{G}_{eff}^{(k)}}\ .
    \end{eqnarray}
    We call this a Gaussian decay.
    
    \item {\bf Brownian ensemble:} In theories with Brownian random couplings, since the Hamiltonian is time-dependent,
    there is no canonical notion of thermal state. Instead, we will focus on the
    infinite temperature state. In this case, we can perform the average exactly
    for all $N$, and the result is
    \begin{eqnarray}
        \mathcal{E}_g\lb U^{\otimes k}\otimes (U^*)^{\otimes k} \rb=e^{-t\mathcal{G}_{eff}^{(k)}}
    \end{eqnarray}
    with the same generator $\mathcal{G}_{eff}^{(k)}$.

    \item We study the generators above in the explicit examples above of fermionic (SYK) and bosonic spin systems. 
\end{enumerate}


    
    \section{To do list:}
    \begin{enumerate}
        \item Generalize the derivation of the $k$-twirl map in the time-independent case at large $N$ to finite temperature.

        \item Understand the physical implications of the $t^{2}$ in the exponent
            of the $k$-twirl map in the large $N$ limit of regular SYK. At late
            time, the decay seems too fast.

        \item Analyze the spectrum of $\mathcal{G}_{eff}^{(k)}$, the invariant subalgebra and establish a spectral gap.

        \item Think about the gluing assumption 1.4 in Nick's notes. Does it
            hold in SYK? Related issue is this: if we couple two Brownian SYK
            models as Swingle does in his 2021 paper, does it look like the Alice-Bob
            protocol for achieving a $k$-twirl map on a combined system $AB$ by applying
            the $k$-twirl maps on $A$ and $B$, separately, swapping $r$ qubits,
            and applying the local $k$-twirls again?

        \item The form of the generators of the Lindbladian is a sum over a hypergraph with hyperedges indexed by $\{i_1,\cdots ,i_q\}$ distinct. go to eq 9 of Nick's notes, and generalize to hypergraphs, with a fermion on each site. 

    Work to bring eq 9 closer to the SYK case by putting fermions on each vertex. 

    Does the gluing lemma hold when you put fermionic degrees of freedom on vertices?

    \item Comment on the $q$-deformed limit $N\to \infty$ with $q^{2}/N$ kept fixed.

    \end{enumerate}


    %\end{document}

    \section{Gaussian decay from time-independent infinite-temperature ensemble}

    Consider a quantum system with Hamiltonian
    \begin{eqnarray}
        H(g)=\sum_I g_I V_I
    \end{eqnarray}
    where $I$ is some collective index, $V_{I}$ are dimensionless self-adjoint operators and $g_{I}$ are real
    coupling parameters with units of energy. We choose $g_{I}$ at random with a
    probability measure that, for simplicity, we take to be Gaussian:
    \begin{eqnarray}
        d\mu(g)= \prod_I\lb dg_I
        \frac{e^{-\frac{1}{2\sigma} g_I^2}}{\sqrt{2\pi \sigma}} \rb
    \end{eqnarray}
    where $\sigma$ has units of energy squared and sets the width of the
    distribution, i.e., $[\sigma]=2$ and $[g_{I}]=1$. The first moments are
    \begin{eqnarray}
        \mathcal{E}_g(g_I)=0, \qquad \mathcal{E}_g(g_I g_J)=\delta_{IJ}\sigma\ .
    \end{eqnarray}

    We represent the thermal state $\rho_{\beta}\propto e^{-\beta H}$ by its canonical
    purification (thermofield double) $\ket{\rho_\beta^{1/2}}$ and the corresponding
    GNS Hilbert space $\mathcal{H}_{\beta}$, generated by the action of
    observables $\mO$ on this vector:
    \begin{eqnarray}
        \mO\to \mO\ket{\rho_\beta^{1/2}}\ .
    \end{eqnarray}
    It is convenient to choose $\beta$ to set the unit of time and define dimensionless
    parameters and Hamiltonian
    \begin{eqnarray}
        &&\tilde{g}_I=\beta g_I,\qquad \tilde{\sigma}=\beta^2 \sigma, \qquad
        \tilde{t}=t/\beta\nn\\
        &&\tilde{H}=\sum_I \tilde{g}_I V_I
    \end{eqnarray}
    The thermal state is a function of $\tilde{g}_{I}$; its canonical
    purification is
    \begin{eqnarray}
        \ket{\rho_\beta^{1/2}}=\ket{\rho_{\tilde{g}}^{1/2}}=(e^{-\tilde{H}/2}\otimes
        1)\ket{1}\ .
    \end{eqnarray}
    where
    \begin{eqnarray}
        \ket{1}\equiv \frac{1}{\sqrt{d}}\sum_E \ket{E E}
    \end{eqnarray}
    and $\{\ket{E}\}$ is the eigenbasis of the Hamiltonian. Since, in general, members
    of our ensemble of Hamiltonians do not commute, we choose some basis for
    $\ket{1}$, and then use the mirror equation
    \begin{eqnarray}
        (\mathcal{O}\otimes 1)\ket{1}=(1\otimes \mathcal{O}^T)\ket{1}
    \end{eqnarray}
    where the transpose is defined in the basis of $\ket{1}$. For self-adjoint
    operators like terms in the Hamiltonians, we have
    \begin{eqnarray}
        (V_I\otimes 1)\ket{1}=(1\otimes V_I^*)\ket{1}\ .
    \end{eqnarray}

    The measure for our ensemble average is
    \begin{eqnarray}
        d\mu(\tilde{g})=\prod_I\lb
        \frac{e^{-\tilde{g}_I^2/(2\tilde{\sigma})}}{\sqrt{2\pi\tilde{\sigma}}} d\tilde{g}_I\rb
    \end{eqnarray}
    For fixed $\tilde{g}$, the time evolution is a unitary super-operator
    \begin{eqnarray}
        \mO(t)=e^{-i H t}\mO e^{i H t}=e^{-i \tilde{H}\tilde{t}} \mO e^{i\tilde{H}\tilde{t}}\ .
    \end{eqnarray}
    It is represented on the Hilbert space as the GNS action of the modular Hamiltonian
    $\tilde{H}$:
    \begin{eqnarray}
        &&\mO(t)\ket{1}=e^{-i \tilde{H} \tilde{t}}\mO\ket{1}\nn\\
        &&\tilde{H}=\sum_I \tilde{g}_I \hat{V}_{I},\qquad \hat{V}_I=V_{I,L}-V^T_{I,R},\qquad
        \hat{V}_I\ket{1}=0.
    \end{eqnarray}
    Here, it is important that $[e^{i \tilde{H} \tilde{t}},e^{-\tilde{H}/2}]=0$. The $g$-averaged time evolution is a unital linear CP map
    \begin{eqnarray}
        \mathcal{E}_g(\mO(t))=\int d\mu(g) e^{-i H(g) t}\mO e^{i H(g) t}\ .
    \end{eqnarray}
    %which as we will show below does not preserve the thermal state $\rho_\beta$!
    In the GNS Hilbert space, the unital CP map $\mathcal{E}_{g}(\cdot)$
    corresponds to the action of the operator
    \begin{eqnarray}
        \mathcal{E}_g(\mO(t))\ket{1}=\int d\mu(g) e^{-i \tilde{t}\sum_I \tilde{g}_I\hat{V}_I }\mO\ket{1}\ .
    \end{eqnarray}

    \subsection{Gaussian average and ordered Wick expansion}

    When there is only a single coupling $g$ (i.e. one index $I$), the Gaussian
    integral can be performed explicitly,
    \begin{eqnarray}
        \int \frac{dg}{\sqrt{2\pi \tilde{\sigma}}}\:e^{-\frac{g^{2}}{2\tilde{\sigma}}}\:e^{-i\tilde{t}g\hat{V}}
        =e^{-\frac{\tilde{\sigma}}{2}\tilde{t}^2\hat{V}^2}\ .
    \end{eqnarray}
    In fact, this integral can be extended to the complex $\tilde{t}$ domain that we denote by $z$ (see appendix \ref{}). Since we are working at infinite temperature, we will not need this analytic continuation for now.
    
    %Therefore,
    % for the $k$-twirl map we have
    % \begin{align}
    %     \int d\mu(g) \mathbb{U}_{\vec{z}}=e^{-\frac{\tilde{\sigma}}{2}\lb\sum_{a=1}^k (iz_a)\hat{V}_{(a)}\rb^2}
    % \end{align}
    % and
    % \begin{align}
    %     G_{1\cdots (k-1)k}(\tau_{1},\cdots,\tau_{k}) & =\braket{\mathbb{S}_{\pi}1^{\otimes k}|e^{-\frac{\tilde{\sigma}}{2}\lb\sum_{a=1}^k (iz_a)\hat{V}_a\rb^2} \otimes_{a=1}^k (\mO_{a})_L|1^{\otimes k}}\nn                   \\
    %                                                  & = \braket{1^{\otimes k}|e^{\frac{\tilde{\sigma}}{2}\lb\sum_{a=1}^k (iz_a)(V_L^{\pi(a)}-(V_R^{(a)})^T)\rb^2} \mathbb{S}_\pi \otimes_{a=1}^k (\mO_{a})_L|1^{\otimes k}}\nn \\
    %                                                  & = \braket{1^{\otimes k}|e^{\frac{\tilde{\sigma}}{2}\lb\sum_{a=1}^k (iz_a)(V_L^{\pi(a)}-(V_R^{(a)})^T)\rb^2} \otimes_{a=1}^k (\mO_{\pi(a)})_L|1^{\otimes k}}
    % \end{align}
    % where we have used
    % \begin{align}
    %     \mathbb{S}_{\pi}(\otimes_{a=1}^{k}(\mO_{a})_{L}) \mathbb{S}_{\pi}^{\dagger}= \otimes_{a=1}^{k}(\mO_{\pi(a)})_{L}.
    % \end{align}
    % We also have
    % \begin{align}
    %     (V_{L}^{(\pi(a))}-(V_{R}^{(a)})^{T})\ket{1^{\otimes k}}= (V_{L}^{(\pi(a))}-V_{L}^{(a)})\ket{1^{\otimes k}}
    % \end{align}
    % Therefore, we can write
    % \begin{align}
    %     G_{1\cdots (k-1)k}(\tau_{1},\cdots,\tau_{k}) & =\braket{1^{\otimes k}|e^{\frac{\tilde{\sigma}}{2}\lb\sum_{a=1}^k (iz_a)(V_L^{\pi(a)}-V_L^{(a)})\rb^2} \otimes_{a=1}^k (\mO_{\pi(a)})_L|1^{\otimes k}}
    % \end{align}

    % It is clear from (\ref{1ptthermal})that for the one-point function individual
    % elements of the ensemble is time-independent. For the $g$-averaged thermal
    % two-point function with insertions at Euclidean times $\tau_{1}$ and
    % $\tau_{2}$ we have
    % \begin{align}
    %     G_{12}(\tau_{1},\tau_{2}) & =\braket{\mathbb{S}_{\pi_0}1^{\otimes 3}| e^{-\frac{\tilde{\sigma}}{2}\lb\tau\hat{V}_2+(1/2-\tau)\hat{V}_3\rb^2} ((\mO_1)_L\otimes (\mO_2)_L\otimes 1_L)|1^{\otimes 3}} %%
    % \end{align}

    % therefore, we have
    % \begin{align}
    %     \braket{\mathbb{S}_{\pi}1^{\otimes 3}|e^{-\frac{\tilde{\sigma}}{2}\lb\tau\hat{V}_2+(1/2-\tau)\hat{V}_3\rb^2} ((\mO_1)_L\otimes (\mO_2)_L\otimes 1_L)|1^{\otimes 3}} & = \braket{1^{\otimes 3}|e^{\frac{\tilde{\sigma}}{2}\lb\tau\hat{V}_2+(1/2-\tau)\hat{V}_3\rb^2} \mathb{S}_\pi ((\mO_1)_L\otimes (\mO_2)_L\otimes 1_L)|1^{\otimes 3}}\nn \\
    %                                                                                                                                                                         & = \braket{1^{\otimes 3}|e^{\frac{\tilde{\sigma}}{2}\lb\tau\hat{V}_2+(1/2-\tau)\hat{V}_3\rb^2} ((\mO_1)_L\otimes (\mO_2)_L\otimes 1_L)|1^{\otimes 3}}
    % \end{align}

    % The $g$-averaged thermal one-point function is
    % \begin{eqnarray}
    %     \braket{\rho_\beta^{1/2}|\mathcal{E}_g\mO(t)\rho_\beta^{1/2}}&=\braket{e^{-\frac{\tilde{\sigma}}{2}\tilde{t}^2\hat{V}^2}\rho_\beta^{1/2}|\mO\rho_\beta^{1/2}}\nn\\
    %     &=
    % \end{eqnarray}
    % which is clearly time-dependent. Hence, $\mathcal{E}_{g}$ does not preserve
    % $\rho_{\beta}$. In particular, in the limit $\tilde{t}\to \infty$, the map projects
    % onto the kernel of $\hat{V}^{2}$. The $g$-averaged two-point function is
    % \begin{eqnarray}
    %     &&\mathcal{E}_g(\braket{\mO(t_1)\rho_\beta^{1/2}|\mO(t_2)\rho_\beta^{1/2}})=\braket{e^{-\frac{\tilde{\sigma}}{2}\hat{V}^2(t_1-t_2)^2}\mO\rho_\beta^{1/2}|\mO\rho_\beta^{1/2}}
    % \end{eqnarray}
    % which is translation-invariant.
    % \NL{When there is only one $g$, for the eigenvectors of $\tilde{V}$ of frequency $\omega$:
    %   \begin{eqnarray}
    %     &&e^{-i g V t}(\mO_\omega)e^{i g V t}=e^{-i\omega g t}\mO_\omega\nn\\
    %     &&\mathcal{E}_g\tr(\rho_\beta \mO_\omega(t))=e^{-\frac{\sigma}{2}(\omega t)^2}\tr(\rho_\beta \mO_\omega)\ .
    %   \end{eqnarray}
    % However, note that these are not modular eigenmodes, because the modular Hamiltonian is $\tilde{H}=g \tilde{V}$ and not just $\tilde{V}$.}

    % \NL{For a single $g$, note that if instead of $g$ we average over modular time with measure
    %   \begin{eqnarray}
    %     d\nu(t)=\sqrt{\frac{\sigma}{2\pi}}e^{-\frac{\sigma t^2}{2}}dt
    %   \end{eqnarray}
    %   we find
    %   \begin{eqnarray}
    %     \int d\nu(t) e^{-i g \tilde{V}t}\mO\ket{\rho^{1/2}}=e^{-\frac{g^2}{2\sigma}\tilde{V}^2}\mO\ket{\rho_\beta^{1/2}}\ .
    %   \end{eqnarray}
    % Does this imply a relation between $g$-average and time average? Should we have started with $H=V$ from the start?}

    More generally, if all the operators $\hat{V}_{I}$ mutually commute, the
    same calculation can be applied in a common eigenbasis and yields
    \begin{eqnarray}
        \int d\mu(g)\:e^{-i\tilde{t}\sum_I g_I\hat{V}_I}=e^{-\frac{\tilde{\sigma}}{2}\tilde{t}^2\sum_I \hat{V}_I^2}\ .
    \end{eqnarray}
When the $\widehat V_I$ mutually commute, the disorder average therefore
produces an exact Gaussian decay.  For noncommuting $\widehat V_I$ the
Gaussian integral does not reduce to the exponential of
$\sum_I\widehat V_I^2$; the deviations are controlled by the operator
ordering in the Wick expansion.

    For noncommuting $\hat{V}_{I}$, the Gaussian integral above does not
    collapse to a single exponential in $\sum_{I}\hat{V}_{I}^{2}$. An exact expression
    can be written as an ordered Wick expansion. Expanding
    $e^{-i\tilde{t}\sum_I g_I\hat{V}_I}$ in a power series and using Wick's
    theorem gives
        \begin{align}
    \int d\mu(g)\,
    e^{-i\tilde{t}\sum_Ig_I\widehat V_I}
    &=
    \sum_{n=0}^{\infty}
    \frac{(-1)^n(\sigma \tilde{t}^2)^n}{(2n)!}
    \sum_{P\in\mathcal P_2(2n)}
    \sum_{I_1,\ldots,I_n}
    \widehat V_{I_{\pi_P(1)}}\cdots
    \widehat V_{I_{\pi_P(2n)}} .
\end{align}
    % \begin{eqnarray}
    %     &&\int d\mu(g)\:e^{-i\tilde{t}\sum_I g_I\hat{V}_I}\nn\\
    %     &&\qquad=\sum_{n=0}^\infty \frac{(-1)^{n}}{(2n)!}\tilde{t}^{2n}\sum_{I_1,\ldots,I_{2n}}\mathcal{E}_g\lb
    %     g_{I_1}\cdots g_{I_{2n}}\rb\;\hat{V}_{I_1}\cdots\hat{V}_{I_{2n}}\nn\\
    %     &&\qquad=\sum_{n=0}^\infty \frac{(-1)^{n}}{2^{n}n!}\tilde{\sigma}^n\tilde{t}^{2n}\sum_{\text{pairings }P}\sum_{I_1,\ldots,I_n}\hat{V}_{I_{\pi_P(1)}}\hat{V}_{I_{\pi_P(2)}}\cdots\hat{V}_{I_{\pi_P(2n)}}\ ,
    % \end{eqnarray}
    and $P$ runs over pairings of $\{1,\ldots,2n\}$ and $\pi_{P}$ is the map
    that repeats each label $I_{r}$ in the two paired slots. The operator ordering
    is inherited from the product inside $\hat{V}_{I_1}\cdots\hat{V}_{I_{2n}}$.
    In particular,
    \begin{eqnarray}
        \int d\mu(g)\:e^{-i\tilde{t}\sum_I g_I\hat{V}_I} &&=1-\frac{\tilde{\sigma}}{2}\tilde{t}^2\sum_I
        \hat{V}_I^2\nn\\
        &&\qquad+\frac{\tilde{\sigma}^{2}}{24}\tilde{t}^4\sum_{I,J}\lb \hat{V}_I^2\hat{V}_J^2+\hat{V}_I\hat{V}_J\hat{V}_I\hat{V}_J+\hat{V}_I\hat{V}_J^2\hat{V}_I\rb+O(\tilde{t}^6)\ ,
    \end{eqnarray}
    Eq.~(17) can also be written as an exponential, but in general it is the exponential
    of an infinite series of nested commutators (a Magnus type expansion).
    Concretely, one may write
    \begin{eqnarray}
        \int d\mu(g)\:e^{-i\tilde{t}\sum_I g_I\hat{V}_I}=\exp\big(\Omega(\tilde{t})\big),\qquad
        \Omega(\tilde{t})=\sum_{n\geq 1}\Omega_n(\tilde{t})\ .
    \end{eqnarray}
    where each $\Omega_{n}$ is built from $n$-fold commutators of the $\hat{V}_{I}$
    and starts at order $\tilde{t}^{2n}$ (odd orders vanish by Gaussianity). In
    general, the $\hat{V}_{I}$ do not commute, $\Omega(\tilde{t})$ does not
    truncate, and we can write the full integral as
    \begin{eqnarray}
        \int d\mu(g)\:e^{-i\tilde{t}\sum_I g_I\hat{V}_I}=\lb \lb e^{\frac{\tilde{\sigma}}{2}\sum_I \frac{\p^{2}}{\p g^{2}_{I}}}\rb
        e^{-i \tilde{t}\sum_I g_I \hat{V}_I}\rb_{g_I=0}\ .
    \end{eqnarray}

    \paragraph{The $k$-twirl channel and its moment operator.}
    For a fixed realization $J$, define $U_{J}(t)=e^{-itH(J)}$ and the $k$-fold
    adjoint action
    \begin{equation}
        \mathcal{U}^{(k)}_{J,t}(Y)\;:=\;U_{J}(t)^{\otimes k}\,Y\,U_{J}(t)^{\dagger\otimes
        k}, \qquad Y\in \mathrm{End}(\mathcal{H}^{\otimes k}).
    \end{equation}
    The disorder-averaged $k$-twirl map is
    \begin{equation}
        \mathcal{E}^{(k)}_{t}(Y)\;:=\;\mathbb{E}_{J}\!\left[\mathcal{U}^{(k)}_{J,t}
        (Y)\right].
    \end{equation}

    Introduce the vectorization/Choi-GNS map $Y\mapsto |Y\rangle\rangle$ on a doubled
    space,
    \begin{equation}
        |Y\rangle\rangle := (Y\otimes \mathbf{1})\,|1\rangle, \qquad |1\rangle :=
        \frac{1}{\sqrt{d}}\sum_{E}|E\rangle\otimes |E\rangle,
    \end{equation}
    so that left- and right-actions are represented by
    \begin{equation}
        (O_{L}- O_{R}^{T})\,|Y\rangle\rangle = |[O,Y]\rangle\rangle.
    \end{equation}
    Then the averaged moment operator (Choi operator) is
    \begin{equation}
        M_{t}^{(k)}\;:=\;\mathbb{E}_{J}\!\left[U_{J}(t)^{\otimes k}\otimes \big(U
        _{J}(t)^{*}\big)^{\otimes k}\right], \qquad |\mathcal{E}_{t}^{(k)}(Y)\rangle
        \rangle = M_{t}^{(k)}\,|Y\rangle\rangle.
    \end{equation}

    %------------------------------------------------------------
    \paragraph{Hatted operators and Gaussian integral.}
    Define, for each replica $a=1,\dots,k$,
    \begin{equation}
        \widehat H_{J}^{(a)}:= H_{J,L}^{(a)}- \big(H_{J,R}^{(a)}\big)^{T}, \qquad
        \widehat H_{J}^{(k)}:= \sum_{a=1}^{k}\widehat H_{J}^{(a)}.
    \end{equation}
    Since $H(J)=\sum_{I}J_{I}V_{I}$, we can write
    \begin{equation}
        \widehat H_{J}^{(k)}\;=\; \sum_{I}J_{I}\,\widehat V_{I}^{(k)}, \qquad \widehat
        V_{I}^{(k)}:= \sum_{a=1}^{k}\widehat V_{I}^{(a)}, \qquad \widehat V_{I}^{(a)}
        := V_{I,L}^{(a)}- \big(V_{I,R}^{(a)}\big)^{T}.
    \end{equation}
    Therefore
    \begin{equation}
        M_{t}^{(k)}\;=\; \mathbb{E}_{J}\!\left[e^{-it\,\widehat H_J^{(k)}}\right]
        \;=\;\int d\mu(J)\;\exp\!\left(-it\sum_{I}J_{I}\,\widehat V_{I}^{(k)}\right
        ), \label{eq:Mt_as_gaussian_integral}
    \end{equation}
    where $d\mu(J)$ is the product Gaussian measure.

    % More generally, we can consider
    % \begin{eqnarray}
    %     M(\tau_1,\cdots,\tau_k)&&:=\mathbb{E_g}\!\left[\otimes_{a=1}^k \lb e^{-\tau_a H_L^{(a)})}\otimes
    %     \big(e^{-\tau_a^* H_R^{(a)}}\big)\rb^{\otimes k}\right],\nn\\
    %     &&=\mathbb{E}_g\left[ e^{-i \sum_I g_I \mathbb{F}_I}\right]\nn\\
    %     \mathbb{F}_I&&=\sum_{a=1}^k(\tau_a V^{(a)}_{I,L}-\tau_a^* V^{(a)}_{I,R})\nn\\
    %     &&\tau_a=\beta_a/4+it_a\ .
    % \end{eqnarray}
    % Note that as a superoperator the operator above for $k=1$ acts as
    % \begin{eqnarray}
    %     &&\mathcal{F}(Y)=e^{-(\beta/4+it)}Y e^{-(\beta/4-it)}\nn\\
    %     &&(\mathcal{F}(Y_L)\otimes 1_R)\ket{e}=(U_t Y_L U_t^\dagger\otimes 1_R)\ket{\rho_\beta^{1/2}}\ .
    % \end{eqnarray}

    %------------------------------------------------------------
    \paragraph{Ordered Wick expansion and leading Gaussian decay.}
    Expanding the exponential and averaging with Wick's theorem gives the exact ordered
    pairing series
    \begin{equation}
        M_{t}^{(k)}= \sum_{n=0}^{\infty}\frac{(-1)^{n}}{2^{n}n!}\,(\sigma t^{2})^{n}
        \sum_{\text{pairings }P}\; \sum_{I_1,\dots,I_n}\widehat V_{I_{\pi_P(1)}}^{(k)}
        \widehat V_{I_{\pi_P(2)}}^{(k)}\cdots \widehat V_{I_{\pi_P(2n)}}^{(k)}, \label{eq:ordered_wick_quenched}
    \end{equation}
    so in particular
    \begin{equation}
        M_{t}^{(k)}= \mathbf{1}-\frac{\sigma t^{2}}{2}\sum_{I}\big(\widehat V_{I}
        ^{(k)}\big)^{2}+\mathcal{O}(t^{4}). \label{eq:Mt_smallt}
    \end{equation}

    Equivalently, one may write a Magnus/cumulant form
    \begin{equation}
        M_{t}^{(k)}=\exp\!\big(\Omega(t)\big), \qquad \Omega(t)=\Omega_{1}(t)+\Omega
        _{2}(t)+\cdots,
    \end{equation}
    where $\Omega_{1}(t)=-(\sigma t^{2}/2)\sum_{I}(\widehat V_{I}^{(k)})^{2}$,
    and $\Omega_{n\ge 2}(t)$ are built from nested commutators and begin at
    order $t^{2n}$.

\subsection{Gaussian decay from large $N$ limit}

Now, consider the system to be comprised of $N$ degrees of freedom,
$i=1,\cdots,N$, and $V_I$ is an interaction term that couples different
sites. For instance, if the sites sit on a lattice, it is natural to consider
nearest-neighbor two-body interactions $V_I$ with $I=\{i,j\}$ that are
adjacent on the lattice. More generally, we can consider the interaction to
be non-local and $q$-body as in the case of hypergraphs such that the label
$I$ is a set of $q$ distinct sites,
\begin{equation}
    I=\{i_1,i_2,\cdots,i_q\}.
\end{equation}
Assuming that the order of labels does not matter, we have the interaction
\begin{equation}
    H
    =
    \sum_{i_1<\cdots<i_q}
    g_{i_1\cdots i_q}V_{i_1\cdots i_q},
    \qquad
    V_{i_1\cdots i_q}
    =
    O_{i_1}\cdots O_{i_q}.
    \label{eq:qbody_general}
\end{equation}
If the operators satisfy $[O_i,O_j]\propto\delta_{ij}$, then
$[V_I,V_J]=0$ for all $I,J$, and we obtain
\begin{equation}
    \int d\mu(g)\,
    e^{-it\sum_I g_I\widehat V_I}
    =
    e^{-\frac{\sigma t^2}{2}
    \sum_I\widehat V_I^2}.
    \label{eq:qbody_commuting_gaussian}
\end{equation}

If there are Majorana fermions sitting at each site, i.e.
\begin{equation}
    \{\psi_i,\psi_j\}=2\delta_{ij},
\end{equation}
there are certain non-vanishing $[V_I,V_J]$. If the two sets of indices
$I$ and $J$ have no overlaps, i.e. $I\cap J=\varnothing$, we have
$[V_I,V_J]=0$. We will assume that $q$ is even. Then, if the number of
overlaps between the two sets,
\begin{equation}
    m=|I\cap J|,
\end{equation}
is even, the interaction terms commute, and if it is odd they anticommute:
\begin{equation}
    [V_I,V_J]
    =
    \left(1-(-1)^{|I\cap J|}\right)V_IV_J.
    \label{eq:majorana_overlap_commutator}
\end{equation}
The case with a Majorana fermion at each site is the SYK model that we will
discuss in the next section.

Here, we simply make the following observation. Choosing random $q$-element
sets $I,J$, the probability of having at least one index in common is
\begin{equation}
    P(\mathrm{overlap})
    =
    1-
    \frac{\binom{N-q}{q}}{\binom Nq}.
    \label{eq:overlap_probability}
\end{equation}
In the limit of fixed $q$ and large $N$,
\begin{equation}
    P(\mathrm{overlap})
    =
    \frac{q^2}{N}
    +
    O(N^{-2})
    \longrightarrow0.
    \label{eq:overlap_largeN}
\end{equation}
Therefore almost all pairs of interaction terms commute at large $N$.
This observation suggests that the Gaussian decay
\begin{equation}
    \int d\mu(g)\,
    e^{-it\sum_I g_I\widehat V_I}
    \simeq
    e^{-\frac{\sigma t^2}{2}
    \sum_I\widehat V_I^2}
    \label{eq:largeN_gaussian_schematic}
\end{equation}
should emerge at large $N$.  The fact that the number of interaction terms
also grows with $N$ means that the vanishing overlap probability by itself
is not sufficient to establish this result. Below, we combine the overlap
counting with the connected cumulant expansion and show that the
non-commuting corrections are parametrically suppressed in the appropriate
large-$N$ scaling regime.


\subsection{Gaussian decay for the $k$-twirl at large $N$}
\label{subsec:quenched_syk_gaussian_regulator}

Let $H(J)$ be a quenched SYK Hamiltonian written as a Gaussian linear
combination
\begin{equation}
    H(J)
    =
    \sum_I J_I W_I,
    \qquad
    W_I:=i^{q/2}V_I=W_I^\dagger,
    \label{eq:quenched_syk_H_general}
\end{equation}
with independent centered Gaussian couplings
\begin{equation}
    \mathbb E[J_I]=0,
    \qquad
    \mathbb E[J_IJ_J]
    =
    \sigma\delta_{IJ},
    \label{eq:quenched_syk_variance}
\end{equation}
and we take $q$ fixed and $N\rightarrow\infty$, with the usual SYK scaling
\begin{equation}
    \sigma\sim N^{-(q-1)}
\end{equation}
understood.

For each replica define
\begin{equation}
    \widehat W_I^{(a)}
    :=
    W_{I,L}^{(a)}
    -
    \left(W_{I,R}^{(a)}\right)^T,
    \qquad
    \widehat W_I^{(k)}
    :=
    \sum_{a=1}^k\widehat W_I^{(a)}.
\end{equation}
Then the moment operator of the $k$-twirl is
\begin{equation}
    M_t^{(k)}
    =
    \mathbb E_J
    \left[
       e^{-it\sum_IJ_I\widehat W_I^{(k)}}
    \right].
    \label{eq:quenched_Mk}
\end{equation}

In models in which the operators $W_I$ are approximately commuting at
large system size---e.g. because most pairs $W_I,W_J$ have disjoint support
and hence commute while commutators require rare overlaps---the
non-commuting corrections in the ordered Wick expansion are suppressed.
To make this statement precise, write
\begin{equation}
    M_t^{(k)}
    =
    e^{\Omega(t)},
    \qquad
    \Omega(t)
    =
    \Omega_1(t)+\Omega_2(t)+\cdots .
\end{equation}
The leading contribution is
\begin{equation}
    \Omega_1(t)
    =
    -\frac{\sigma t^2}{2}
    \sum_I
    \left(\widehat W_I^{(k)}\right)^2.
    \label{eq:quenched_Omega1}
\end{equation}

At fourth order, using the Wick expansion of Sec.~2.1, one finds
\begin{align}
    \Omega_2(t)
    =
    \frac{\sigma^2t^4}{24}
    \sum_{I,J}
    \bigg[
       &
       \widehat W_I^{(k)}
       \widehat W_J^{(k)}
       \widehat W_I^{(k)}
       \widehat W_J^{(k)}
       +
       \widehat W_I^{(k)}
       \left(\widehat W_J^{(k)}\right)^2
       \widehat W_I^{(k)}
       \nonumber\\
       &
       -
       2
       \left(\widehat W_I^{(k)}\right)^2
       \left(\widehat W_J^{(k)}\right)^2
    \bigg].
    \label{eq:quenched_Omega2}
\end{align}
The expression in square brackets vanishes whenever
\begin{equation}
    [\widehat W_I^{(k)},\widehat W_J^{(k)}]=0.
\end{equation}
Thus $\Omega_2$ receives contributions only from the $O(1/N)$ fraction of
pairs whose supports overlap in such a way that the corresponding
interaction terms fail to commute.

There are
\begin{equation}
    M_N=\binom Nq
\end{equation}
interaction terms in SYK.  We consider the diffusive scaling regime
\begin{equation}
    \tau
    :=
    \sigma t^2\binom Nq
    \sim Nt^2
    =
    O(1).
    \label{eq:quenched_diffusive_tau}
\end{equation}
Then
\begin{equation}
    \Omega_1=O(1),
\end{equation}
whereas
\begin{equation}
    \Omega_2
    =
    O\left(
       \sigma^2t^4
       \binom Nq^2
       \frac1N
    \right)
    =
    O(N^{-1}).
    \label{eq:quenched_Omega2_scaling}
\end{equation}

The same counting extends order by order in the logarithm of the moment
operator.  Consider a contribution to $\Omega_n$ containing $r\leq n$
distinct interaction labels.  Construct a graph whose vertices are the
distinct interaction terms and join two vertices when the corresponding
operators fail to commute.  If this graph is disconnected, the corresponding
sets of operators commute with one another and the Gaussian average
factorizes between the connected components.  Such a disconnected
contribution therefore does not appear in the logarithm.

A contribution to $\Omega_n$ must consequently have a connected
non-commutation graph.  For fixed $q$ and fixed $r$, after choosing the
first interaction freely, every additional vertex of a spanning tree must
overlap one of the previously chosen interactions.  Each such overlap costs
a factor $O(1/N)$.  Hence the number of connected $r$-tuples scales as
\begin{equation}
    O\left(
       \frac{M_N^r}{N^{r-1}}
    \right).
\end{equation}
It follows that, at fixed order $n$,
\begin{align}
    \Omega_n
    &=
    O\left[
       (\sigma t^2)^n
       \frac{M_N^r}{N^{r-1}}
    \right]
    \nonumber\\
    &=
    O(N^{-(n-1)}),
    \qquad
    n\geq2,
    \label{eq:quenched_Omegan_scaling}
\end{align}
where we have used $M_N=\Theta(N^q)$ and
$\sigma=\Theta(N^{-(q-1)})$ together with
$\tau=\sigma M_Nt^2=O(1)$.

Therefore, order by order in the fixed-$q$, fixed-$k$ large-$N$ expansion,
\begin{equation}
    M_t^{(k)}
    \simeq
    \exp\left[
       -\frac{\sigma t^2}{2}
       \sum_I
       \left(\widehat W_I^{(k)}\right)^2
    \right].
    \label{eq:quenched_gaussian_regulator}
\end{equation}
Here $\simeq$ denotes equality to leading order at large $N$.  The argument
above establishes the suppression of the non-Gaussian terms at every fixed
order in $\tau$.  A uniform estimate on the fully resummed series at
arbitrary finite $\tau$ would require an additional summability bound.

\paragraph{Identification of the effective generator.}
Define
\begin{equation}
    \widehat X_I
    :=
    \sqrt{\sigma}\,
    \widehat W_I^{(k)}
    =
    \sqrt{\sigma}
    \sum_{a=1}^k
    \left(
       W_{I,L}^{(a)}
       -
       (W_{I,R}^{(a)})^T
    \right),
    \qquad
    \widehat G_{\rm eff}^{(k)}
    :=
    \frac12\sum_I\widehat X_I^2.
    \label{eq:hatGeff_quenched}
\end{equation}
Then Eq.~\eqref{eq:quenched_gaussian_regulator} becomes
\begin{equation}
    \boxed{
    M_t^{(k)}
    \simeq
    e^{-t^2\widehat G_{\rm eff}^{(k)}}.
    }
    \label{eq:Mt_exp_hatGeff}
\end{equation}

On the physical $k$-replica space define Hermitian operators
\begin{equation}
    X_I
    :=
    \sqrt{\sigma}
    \sum_{a=1}^kW_I^{(a)},
    \qquad
    \widehat X_I
    =
    (X_I)_L-(X_I)_R^T.
\end{equation}
Using
\begin{equation}
    \left((X)_L-(X)_R^T\right)^2|Y\rangle\rangle
    =
    |X^2Y+YX^2-2XYX\rangle\rangle,
\end{equation}
we obtain the corresponding superoperator $G_{\rm eff}^{(k)}$ on $Y$ via
\begin{equation}
    \widehat G_{\rm eff}^{(k)}|Y\rangle\rangle
    =
    |G_{\rm eff}^{(k)}(Y)\rangle\rangle,
\end{equation}
namely
\begin{align}
    G_{\rm eff}^{(k)}(Y)
    &=
    \frac12\sum_I
    \left(
       X_I^2Y+YX_I^2-2X_IYX_I
    \right)
    \nonumber\\
    &=
    \frac12\sum_I[X_I,[X_I,Y]]
    \nonumber\\
    &=
    -\sum_I
    \left(
       X_IYX_I
       -
       \frac12\{X_I^2,Y\}
    \right).
    \label{eq:Geff_quenched}
\end{align}
Therefore, in the same large-$N$ regime in which
Eq.~\eqref{eq:Mt_exp_hatGeff} holds,
\begin{equation}
    \boxed{
    {\cal E}_t^{(k)}
    \simeq
    e^{-t^2G_{\rm eff}^{(k)}}.
    }
    \label{eq:quenched_channel_final}
\end{equation}
\section{Gaussian decay in large $N$ bosonic and Spin SYK}
\label{sec:bosonic_spin_syk_gaussian_regulator}

In Sec.~\ref{subsec:quenched_syk_gaussian_regulator} we argued that the
quenched $k$-twirl becomes a Gaussian decay at large $N$ when the
non-commuting interaction terms require rare overlaps of their supports.
The argument applies more generally to fixed-$q$ all-to-all models for
which
\begin{equation}
    M_N=\Theta(N^q),
    \qquad
    \sigma_N=\Theta(N^{-(q-1)}),
\end{equation}
and the probability that two randomly chosen interaction terms fail to
commute is $O(N^{-1})$.  In this section we apply this result to the
bosonic Spin SYK models of \cite{SwingleWiner} and the two-component
Spin SYK models of \cite{BasuDasNandy}.  The overlap counting is
particularly simple in these models because operators acting on different
physical sites commute exactly.

For a Hamiltonian
\begin{equation}
    H(J)=\sum_{I\in{\cal A}_N}J_IV_I,
    \qquad
    \mathbb E[J_I]=0,
    \qquad
    \mathbb E[J_IJ_J]=\sigma_N\delta_{IJ},
    \label{eq:spin_syk_general_H}
\end{equation}
define
\begin{equation}
    W_I
    :=
    \sum_{a=1}^kV_I^{(a)}.
\end{equation}
Whenever the conditions above hold, the result of
Sec.~\ref{subsec:quenched_syk_gaussian_regulator} gives, order by order
in the fixed-$q$, fixed-$k$ large-$N$ expansion,
\begin{equation}
    {\cal E}_t^{(k)}
    \simeq
    e^{-t^2G_{\rm eff}^{(k)}},
    \label{eq:spin_syk_general_regulator}
\end{equation}
where
\begin{equation}
    G_{\rm eff}^{(k)}(Y)
    =
    \frac{\sigma_N}{2}
    \sum_{I\in{\cal A}_N}
    [W_I,[W_I,Y]].
    \label{eq:spin_syk_Geff_general}
\end{equation}
The natural large-$N$ scaling variable is
\begin{equation}
    \tau_N
    :=
    \sigma_NM_Nt^2
    =
    O(1),
    \qquad
    t^2=O(N^{-1}).
    \label{eq:spin_syk_diffusive_scaling}
\end{equation}


\subsection{Three-component bosonic Spin SYK}
\label{subsec:three_component_spin_syk_regulator}

We first consider the Spin SYK model of \cite{SwingleWiner}.  There are
$N$ qubits and a generic interaction is
\begin{equation}
    V_I
    =
    \sigma_{r_1}^{\mu_1}\cdots\sigma_{r_q}^{\mu_q},
    \qquad
    r_1<\cdots<r_q,
    \qquad
    \mu_j\in\{x,y,z\}.
    \label{eq:SW_VI}
\end{equation}
The Hamiltonian is
\begin{equation}
    H
    =
    \sum_{\substack{r_1<\cdots<r_q\\
                    \mu_1,\ldots,\mu_q}}
    J_{r_1\mu_1\cdots r_q\mu_q}
    \sigma_{r_1}^{\mu_1}\cdots\sigma_{r_q}^{\mu_q},
    \label{eq:SW_H}
\end{equation}
with
\begin{equation}
    \mathbb E[J_I]=0,
    \qquad
    \mathbb E[J_IJ_J]
    =
    \sigma_q^{(3)}\delta_{IJ},
    \qquad
    \sigma_q^{(3)}
    =
    \frac{(q-1)!J^2}{N^{q-1}}.
    \label{eq:SW_sigma}
\end{equation}
There are
\begin{equation}
    M_q^{(3)}
    =
    3^q\binom Nq
    \label{eq:SW_number_terms}
\end{equation}
interaction terms.

The overlap counting can be sharpened because two overlapping Pauli
strings need not anticommute.  Suppose two interaction terms overlap on
$m$ sites.  At each shared site the two Pauli matrices disagree with
probability $2/3$.  If $d$ is the number of shared sites with different
Pauli components, then
\begin{equation}
    V_IV_J=(-1)^dV_JV_I.
    \label{eq:SW_commutation_sign}
\end{equation}
Conditioned on an overlap of size $m$,
\begin{equation}
    {\mathbb P}
    \left(
       [V_I,V_J]\neq0
       \,\middle|\,
       |I\cap J|=m
    \right)
    =
    {\mathbb P}(d\ {\rm odd})
    =
    \frac{1-(-1/3)^m}{2}.
    \label{eq:SW_noncommuting_given_m}
\end{equation}
Hence the exact probability that two independently chosen interaction
strings fail to commute is
\begin{equation}
    p_{\rm nc}^{(3)}
    =
    \sum_{m=1}^q
    \frac{
       \binom qm\binom{N-q}{q-m}
    }{
       \binom Nq
    }
    \frac{1-(-1/3)^m}{2}.
    \label{eq:SW_pnc_exact}
\end{equation}
At fixed $q$ and large $N$, the $m=1$ term dominates, giving
\begin{equation}
    p_{\rm nc}^{(3)}
    =
    \frac{2q^2}{3N}
    +
    O(N^{-2}).
    \label{eq:SW_pnc_largeN}
\end{equation}
Thus the additional spin-component labels change only the numerical
coefficient of the rare-overlap probability; the $1/N$ suppression is
unchanged.

The effective $k$-twirl generator is
\begin{equation}
    G_{{\rm eff},q}^{(k,3)}(Y)
    =
    \frac{(q-1)!J^2}{2N^{q-1}}
    \sum_{\substack{r_1<\cdots<r_q\\
                    \mu_1,\ldots,\mu_q}}
    [W_I,[W_I,Y]],
    \label{eq:SW_Geff_q}
\end{equation}
where
\begin{equation}
    W_I
    :=
    \sum_{a=1}^k
    \left(
       \sigma_{r_1}^{\mu_1}
       \cdots
       \sigma_{r_q}^{\mu_q}
    \right)^{(a)}.
    \label{eq:SW_WI}
\end{equation}
The diffusive variable is
\begin{align}
    \tau_q^{(3)}
    &:=
    \sigma_q^{(3)}M_q^{(3)}t^2
    \nonumber\\
    &=
    \frac{(q-1)!J^2}{N^{q-1}}
    3^q\binom Nq t^2
    \nonumber\\
    &=
    \frac{3^q}{q}J^2Nt^2
    \left[
       1+O(N^{-1})
    \right].
    \label{eq:SW_tau_general}
\end{align}

For the quadratic model,
\begin{equation}
    q=2:
    \qquad
    \sigma_2^{(3)}=\frac{J^2}{N},
    \qquad
    M_2^{(3)}=9\binom N2,
    \qquad
    p_{\rm nc}^{(3)}
    =
    \frac{8}{3N}+O(N^{-2}),
    \label{eq:SW_q2_data}
\end{equation}
and
\begin{equation}
    \tau_2^{(3)}
    =
    \frac92J^2Nt^2
    \left[
       1+O(N^{-1})
    \right].
    \label{eq:SW_q2_tau}
\end{equation}
The generator is
\begin{equation}
    G_{{\rm eff},2}^{(k,3)}(Y)
    =
    \frac{J^2}{2N}
    \sum_{r<s}
    \sum_{\mu,\nu=x,y,z}
    [W_{rs}^{\mu\nu},
     [W_{rs}^{\mu\nu},Y]],
    \label{eq:SW_q2_Geff}
\end{equation}
where
\begin{equation}
    W_{rs}^{\mu\nu}
    :=
    \sum_{a=1}^k
    (\sigma_r^\mu\sigma_s^\nu)^{(a)}.
\end{equation}

For the quartic model,
\begin{equation}
    q=4:
    \qquad
    \sigma_4^{(3)}
    =
    \frac{6J^2}{N^3},
    \qquad
    M_4^{(3)}
    =
    81\binom N4,
    \qquad
    p_{\rm nc}^{(3)}
    =
    \frac{32}{3N}
    +
    O(N^{-2}),
    \label{eq:SW_q4_data}
\end{equation}
and
\begin{equation}
    \tau_4^{(3)}
    =
    \frac{81}{4}J^2Nt^2
    \left[
       1+O(N^{-1})
    \right].
    \label{eq:SW_q4_tau}
\end{equation}
The generator is
\begin{equation}
    G_{{\rm eff},4}^{(k,3)}(Y)
    =
    \frac{3J^2}{N^3}
    \sum_{r_1<r_2<r_3<r_4}
    \sum_{\mu_1,\ldots,\mu_4=x,y,z}
    [W_I,[W_I,Y]].
    \label{eq:SW_q4_Geff}
\end{equation}
Thus, for $q=2$ and $q=4$,
\begin{equation}
    {\cal E}_{t,q}^{(k)}
    \simeq
    e^{-t^2G_{{\rm eff},q}^{(k,3)}},
    \qquad
    t^2=O(N^{-1}).
    \label{eq:SW_final_regulator}
\end{equation}

This statement concerns the quenched disorder average of the real-time
$k$-twirl in the diffusive large-$N$ limit.  It should not be confused
with the low-temperature saddle-point question studied in
\cite{SwingleWiner}; in particular, the Gaussian-decay counting does not
by itself determine whether the low-energy phase is glassy or SYK-like.


\subsection{Two-component Spin SYK}
\label{subsec:two_component_spin_syk_regulator}

We next consider the Spin SYK models of \cite{BasuDasNandy}.  We denote
the number of physical spin sites by $N_{\rm Spin}$.  The
$2N_{\rm Spin}$ constituent operators are
\begin{equation}
    O_{2r-1}=\sigma_r^x,
    \qquad
    O_{2r}=\sigma_r^y,
    \qquad
    r=1,\ldots,N_{\rm Spin},
    \label{eq:BDN_constituent_O}
\end{equation}
where identities on all other sites are implicit.  The Hamiltonian is
\begin{equation}
    H_{{\rm Spin\,SYK}_q}
    =
    \sqrt{
       \frac{(q-1)!}{(2N_{\rm Spin})^{q-1}}
    }
    \sum_{1\leq i_1<\cdots<i_q\leq2N_{\rm Spin}}
    J_{i_1\cdots i_q}
    V_{i_1\cdots i_q},
    \label{eq:BDN_H}
\end{equation}
with
\begin{equation}
    V_{i_1\cdots i_q}
    :=
    i^{\eta_{i_1\cdots i_q}}
    O_{i_1}\cdots O_{i_q},
    \qquad
    V_{i_1\cdots i_q}
    =
    V_{i_1\cdots i_q}^\dagger,
    \label{eq:BDN_VI}
\end{equation}
and the $J_{i_1\cdots i_q}$ are independent unit-variance Gaussian
random variables.  Equivalently, after absorbing the deterministic
prefactor into the random coupling,
\begin{equation}
    \sigma_q^{(2)}
    =
    \frac{(q-1)!}
         {(2N_{\rm Spin})^{q-1}}.
    \label{eq:BDN_sigma}
\end{equation}

It is useful first to consider the genuine model, denoted gSpin-SYK in
\cite{BasuDasNandy}, in which no interaction contains two constituent
operators from the same physical spin site.  Since all factors then act
on distinct sites, a genuine interaction is simply
\begin{equation}
    V_{S,\boldsymbol\alpha}
    =
    \prod_{r\in S}\sigma_r^{\alpha_r},
    \qquad
    |S|=q,
    \qquad
    \alpha_r\in\{x,y\}.
    \label{eq:BDN_genuine_V}
\end{equation}
There are
\begin{equation}
    M_{q,g}^{(2)}
    =
    2^q
    \binom{N_{\rm Spin}}q
    \label{eq:BDN_genuine_number}
\end{equation}
such terms.

Suppose two genuine interactions overlap on $m\geq1$ physical sites.
On each shared site the two Pauli components disagree with probability
$1/2$.  The number of mismatches is therefore binomially distributed
with parameter $1/2$, and for every $m\geq1$ its parity is exactly
uniform.  Hence
\begin{equation}
    {\mathbb P}
    \left(
       [V_I,V_J]\neq0
       \,\middle|\,
       |S_I\cap S_J|=m
    \right)
    =
    \frac12.
    \label{eq:BDN_conditional_pnc}
\end{equation}
It follows that
\begin{equation}
    p_{{\rm nc},g}^{(2)}
    =
    \frac12
    \left[
       1-
       \frac{
          \binom{N_{\rm Spin}-q}{q}
       }{
          \binom{N_{\rm Spin}}q
       }
    \right]
    =
    \frac{q^2}{2N_{\rm Spin}}
    +
    O(N_{\rm Spin}^{-2}).
    \label{eq:BDN_pnc}
\end{equation}

The effective generator of the genuine model is
\begin{equation}
    G_{{\rm eff},q,g}^{(k,2)}(Y)
    =
    \frac{(q-1)!}
         {2(2N_{\rm Spin})^{q-1}}
    \sum_{\substack{|S|=q\\
                    \boldsymbol\alpha\in\{x,y\}^q}}
    [W_{S,\boldsymbol\alpha},
     [W_{S,\boldsymbol\alpha},Y]],
    \label{eq:BDN_genuine_Geff}
\end{equation}
where
\begin{equation}
    W_{S,\boldsymbol\alpha}
    :=
    \sum_{a=1}^k
    \left(
       \prod_{r\in S}\sigma_r^{\alpha_r}
    \right)^{(a)}.
    \label{eq:BDN_genuine_W}
\end{equation}
Its diffusive variable is
\begin{align}
    \tau_{q,g}^{(2)}
    &:=
    \sigma_q^{(2)}M_{q,g}^{(2)}t^2
    \nonumber\\
    &=
    \frac{(q-1)!}
         {(2N_{\rm Spin})^{q-1}}
    2^q
    \binom{N_{\rm Spin}}q
    t^2
    \nonumber\\
    &=
    \frac{2}{q}
    N_{\rm Spin}t^2
    \left[
       1+O(N_{\rm Spin}^{-1})
    \right].
    \label{eq:BDN_tau_general}
\end{align}

For $q=2$,
\begin{equation}
    \sigma_2^{(2)}
    =
    \frac{1}{2N_{\rm Spin}},
    \qquad
    M_{2,g}^{(2)}
    =
    4\binom{N_{\rm Spin}}2,
    \qquad
    p_{{\rm nc},g}^{(2)}
    =
    \frac{2}{N_{\rm Spin}}
    +
    O(N_{\rm Spin}^{-2}),
    \label{eq:BDN_q2_data}
\end{equation}
and
\begin{equation}
    \tau_{2,g}^{(2)}
    =
    N_{\rm Spin}t^2
    \left[
       1+O(N_{\rm Spin}^{-1})
    \right].
    \label{eq:BDN_q2_tau}
\end{equation}
The generator is
\begin{equation}
    G_{{\rm eff},2,g}^{(k,2)}(Y)
    =
    \frac{1}{4N_{\rm Spin}}
    \sum_{r<s}
    \sum_{\alpha,\beta=x,y}
    [W_{rs}^{\alpha\beta},
     [W_{rs}^{\alpha\beta},Y]],
    \label{eq:BDN_q2_Geff}
\end{equation}
where
\begin{equation}
    W_{rs}^{\alpha\beta}
    =
    \sum_{a=1}^k
    (\sigma_r^\alpha\sigma_s^\beta)^{(a)}.
\end{equation}

For $q=4$,
\begin{equation}
    \sigma_4^{(2)}
    =
    \frac{3}{4N_{\rm Spin}^3},
    \qquad
    M_{4,g}^{(2)}
    =
    16\binom{N_{\rm Spin}}4,
    \qquad
    p_{{\rm nc},g}^{(2)}
    =
    \frac{8}{N_{\rm Spin}}
    +
    O(N_{\rm Spin}^{-2}),
    \label{eq:BDN_q4_data}
\end{equation}
and
\begin{equation}
    \tau_{4,g}^{(2)}
    =
    \frac12N_{\rm Spin}t^2
    \left[
       1+O(N_{\rm Spin}^{-1})
    \right].
    \label{eq:BDN_q4_tau}
\end{equation}
The corresponding generator is
\begin{equation}
    G_{{\rm eff},4,g}^{(k,2)}(Y)
    =
    \frac{3}{8N_{\rm Spin}^3}
    \sum_{r_1<r_2<r_3<r_4}
    \sum_{\alpha_1,\ldots,\alpha_4=x,y}
    [W_I,[W_I,Y]].
    \label{eq:BDN_q4_Geff}
\end{equation}
Thus, for both $q=2$ and $q=4$,
\begin{equation}
    {\cal E}_{t,q,g}^{(k)}
    \simeq
    e^{-t^2G_{{\rm eff},q,g}^{(k,2)}},
    \qquad
    t^2=O(N_{\rm Spin}^{-1}).
    \label{eq:BDN_genuine_final}
\end{equation}


\subsection{Full versus genuine Spin SYK}
\label{subsec:full_vs_genuine_spin_syk}

The full Spin-SYK model of \cite{BasuDasNandy} allows several of the
constituent indices $i_1,\ldots,i_q$ to correspond to different Pauli
components on the same physical site.  Because only $\sigma^x$ and
$\sigma^y$ occur, at most two constituent indices can come from one
site.  The total number of terms in the full model is
\begin{equation}
    M_{q,{\rm full}}^{(2)}
    =
    \binom{2N_{\rm Spin}}q,
    \label{eq:BDN_full_number}
\end{equation}
whereas the number of genuine $q$-site terms is
$2^q\binom{N_{\rm Spin}}q$.  Their ratio obeys
\begin{align}
    \frac{
       M_{q,g}^{(2)}
    }{
       M_{q,{\rm full}}^{(2)}
    }
    &=
    \frac{
       2^q\binom{N_{\rm Spin}}q
    }{
       \binom{2N_{\rm Spin}}q
    }
    \nonumber\\
    &=
    1-
    \frac{q(q-1)}
         {4N_{\rm Spin}}
    +
    O(N_{\rm Spin}^{-2}).
    \label{eq:BDN_genuine_fraction}
\end{align}
Thus the number of collision terms is
\begin{equation}
    M_{q,{\rm coll}}^{(2)}
    =
    O(N_{\rm Spin}^{q-1}).
    \label{eq:BDN_collision_number}
\end{equation}
Since
\begin{equation}
    \sigma_q^{(2)}
    =
    O(N_{\rm Spin}^{-(q-1)}),
\end{equation}
their contribution to the unscaled effective generator is
\begin{equation}
    \Delta G_{\rm eff}^{(k)}
    =
    O(1),
    \label{eq:BDN_collision_G}
\end{equation}
whereas the genuine $q$-site sector is $O(N_{\rm Spin})$.  On the
diffusive time scale $t^2=O(N_{\rm Spin}^{-1})$,
\begin{equation}
    t^2\Delta G_{\rm eff}^{(k)}
    =
    O(N_{\rm Spin}^{-1}).
    \label{eq:BDN_collision_suppression}
\end{equation}
Consequently, for fixed even $q$,
\begin{equation}
    {\cal E}_{t,q,{\rm full}}^{(k)}
    \simeq
    {\cal E}_{t,q,g}^{(k)}
    \simeq
    e^{-t^2G_{{\rm eff},q,g}^{(k,2)}}.
    \label{eq:BDN_full_genuine_same}
\end{equation}
In particular, the full and genuine quadratic models have the same
leading large-$N_{\rm Spin}$ Gaussian decay, and the same is true of the
full and genuine quartic models.  The self-site terms in the quadratic
model and the lower-support collision terms in the quartic model first
affect the diffusive evolution at order $1/N_{\rm Spin}$.


\subsection{Conserved $\mathbb Z_2$ charge in the two-component model}
\label{subsec:BDN_Z2_charge}

There is one additional difference between the two-component model and
the generic three-component bosonic Spin SYK model.  Define
\begin{equation}
    \Gamma^3
    :=
    \bigotimes_{r=1}^{N_{\rm Spin}}\sigma_r^z.
    \label{eq:BDN_Gamma3}
\end{equation}
Every constituent operator in Eq.~\eqref{eq:BDN_constituent_O}
anticommutes with $\Gamma^3$,
\begin{equation}
    \{\Gamma^3,O_i\}=0.
    \label{eq:BDN_Gamma_O}
\end{equation}
It follows immediately that an interaction containing $q$ constituent
operators satisfies
\begin{equation}
    \Gamma^3V_I
    =
    (-1)^qV_I\Gamma^3.
    \label{eq:BDN_Gamma_V}
\end{equation}
Therefore, for the quadratic and quartic models,
\begin{equation}
    [\Gamma^3,V_I]=0,
    \qquad
    q=2,4.
    \label{eq:BDN_even_q_charge}
\end{equation}
The same statement holds for both the full and genuine ensembles.

On the $k$-replica Hilbert space, each replica charge
$\Gamma^{3(a)}$ commutes with all the operators
\begin{equation}
    X_I
    =
    \sqrt{\sigma_q^{(2)}}
    \sum_{b=1}^kV_I^{(b)}.
\end{equation}
Hence
\begin{equation}
    \Gamma^{3(a)}
    \in
    {\rm Comm}\{X_I\},
    \qquad
    a=1,\ldots,k.
    \label{eq:BDN_replica_charge_commutant}
\end{equation}
Together with the replica permutations, these operators therefore lie in
the kernel of $G_{\rm eff}^{(k)}$.  In particular, the two-component
$k$-twirl cannot converge to the unrestricted Haar twirl on the full
Hilbert space.  Determining whether these charges and the replica
permutations generate the complete invariant algebra requires a separate
analysis of the interaction algebra and is not implied by the overlap
counting above.


\subsection{Summary}

The quadratic and quartic bosonic Spin SYK models obey the same large-$N$
mechanism established in Sec.~\ref{subsec:quenched_syk_gaussian_regulator}.
Their interaction terms have bounded support, the number of leading
$q$-body terms is $\Theta(N^q)$, and the probability that two randomly
chosen interaction terms fail to commute is $O(1/N)$.  With the usual
extensive normalization,
\begin{equation}
    \sigma_NM_N=\Theta(N),
\end{equation}
the finite large-$N$ time variable is
\begin{equation}
    \tau_N
    =
    \sigma_NM_Nt^2
    =
    O(1),
    \qquad
    t=O(N^{-1/2}).
\end{equation}
The leading $k$-twirl is therefore
\begin{equation}
    \boxed{
    {\cal E}_t^{(k)}
    \simeq
    e^{-t^2G_{\rm eff}^{(k)}},
    \qquad
    G_{\rm eff}^{(k)}(Y)
    =
    \frac{\sigma_N}{2}
    \sum_I
    \left[
       \sum_{a=1}^kV_I^{(a)},
       \left[
          \sum_{b=1}^kV_I^{(b)},Y
       \right]
    \right].
    }
    \label{eq:bosonic_spin_syk_final_box}
\end{equation}
For the two-component Spin SYK model, the full and genuine ensembles give
the same leading generator, while the exact $\Gamma^3$ charge enlarges
the invariant algebra for even $q$.



\section{Gaussian decay in large $N$ SYK}
\label{sec:largeN_SYK}

We now return to fermionic SYK and study the structure of the effective
generator in more detail.  The SYK model is comprised of $N$ Majorana
fermions $\chi_i$, $i=1,\ldots,N$, with
\begin{equation}
    \{\chi_i,\chi_j\}=2\delta_{ij},
\end{equation}
and Hamiltonian
\begin{equation}
    H
    =
    i^{q/2}
    \sum_{i_1<\cdots<i_q}
    J_{i_1\cdots i_q}
    \chi_{i_1}\cdots\chi_{i_q}.
    \label{eq:SYK_H_static}
\end{equation}
We restrict to even $q$ and even $N$.  In the notation of
Sec.~\ref{subsec:quenched_syk_gaussian_regulator}, define
\begin{equation}
    V_I
    :=
    \chi_{i_1}\cdots\chi_{i_q},
    \qquad
    W_I
    :=
    i^{q/2}V_I.
    \label{eq:SYK_VW}
\end{equation}
It follows from the Majorana algebra that
\begin{equation}
    V_I^2
    =
    (-1)^{q(q-1)/2}
    =
    (-1)^{q/2},
\end{equation}
so that
\begin{equation}
    W_I=W_I^\dagger,
    \qquad
    W_I^2=1.
    \label{eq:SYK_W_square}
\end{equation}

The time-independent noise is
\begin{equation}
    \mathbb E[J_I]=0,
    \qquad
    \mathbb E[J_IJ_J]
    =
    \sigma\delta_{IJ},
\end{equation}
with the normalization used in [1,2],
\begin{equation}
    \sigma
    =
    \frac{2^{q-1}q!}{q^2N^{q-1}}J.
    \label{eq:SYK_sigma}
\end{equation}
There are $\binom Nq$ terms in the Hamiltonian, which for fixed $q$ and
large $N$ grows as $N^q/q!$.  For $q=2$ and $q=4$ this normalization
gives, respectively,
\begin{equation}
    \sigma=\frac{J}{N},
    \qquad
    \sigma=\frac{12J}{N^3}.
\end{equation}


\subsection{Jordan-Wigner transform}

Assuming that $N$ is even, we pair the Majorana fermions
$\chi_{2j-1}$ and $\chi_{2j}$ into $N/2$ complex fermions
\begin{equation}
    c_j
    :=
    \frac12
    \left(
       \chi_{2j-1}
       +
       i\chi_{2j}
    \right),
\end{equation}
which satisfy
\begin{equation}
    \{c_l,c_j^\dagger\}=\delta_{lj},
    \qquad
    \{c_l,c_j\}
    =
    \{c_l^\dagger,c_j^\dagger\}
    =
    0.
\end{equation}
We use the standard qubit convention
\begin{equation}
    \sigma_j^-
    =
    \frac12(X_j+iY_j)
    =
    |0\rangle\langle1|_j,
    \qquad
    N_j
    :=
    c_j^\dagger c_j
    =
    \frac{1-Z_j}{2}.
    \label{eq:SYK_qubit_convention}
\end{equation}
Then
\begin{equation}
    Z_j
    =
    -i\chi_{2j-1}\chi_{2j},
    \qquad
    (-1)^{N_j}=Z_j,
\end{equation}
and the Jordan-Wigner transformation is
\begin{align}
    c_j
    &=
    (Z_1\cdots Z_{j-1})\sigma_j^-,
    \nonumber\\
    \chi_{2j-1}
    &=
    (Z_1\cdots Z_{j-1})X_j,
    \nonumber\\
    \chi_{2j}
    &=
    (Z_1\cdots Z_{j-1})Y_j.
    \label{eq:SYK_JW}
\end{align}

Let
\begin{equation}
    |e\rangle
    =
    \frac{1}{\sqrt d}
    \sum_{\mathbf n}
    |\mathbf n\rangle_L|\mathbf n\rangle_R
\end{equation}
be the maximally entangled vector in this fixed Jordan-Wigner basis.
For any operator $O$,
\begin{equation}
    (O_L-O_R^T)|e\rangle=0.
    \label{eq:SYK_general_mirror}
\end{equation}
It is convenient to define the Majorana operators on the $R$ copy by
\begin{equation}
    \chi_{j,R}
    :=
    (\chi_j^T)_R.
    \label{eq:SYK_R_majorana}
\end{equation}
Then
\begin{equation}
    (\chi_{j,L}-\chi_{j,R})|e\rangle=0,
\end{equation}
while the $L$ and $R$ Majoranas still commute because they act on
different tensor factors.

For the Hermitian SYK interaction
\begin{equation}
    W_I
    =
    i^{q/2}
    \chi_{i_1}\cdots\chi_{i_q},
\end{equation}
transpose reverses the order of the Majoranas.  Therefore
\begin{align}
    (W_I^T)_R
    &=
    i^{q/2}
    \chi_{i_q,R}\cdots\chi_{i_1,R}
    \nonumber\\
    &=
    (-i)^{q/2}
    \chi_{i_1,R}\cdots\chi_{i_q,R}.
    \label{eq:SYK_W_transpose}
\end{align}


\subsection{Totally anti-commuting variables}

Since we are interested in the $k$-twirl moment operator
\begin{equation}
    U^{\otimes k}\otimes(U^*)^{\otimes k},
\end{equation}
consider $k$ copies of the $L$ Majoranas and $k$ copies of the $R$
Majoranas,
\begin{equation}
    \chi_{l,L}^{(a)},
    \qquad
    \chi_{l,R}^{(a)},
    \qquad
    l=1,\ldots,N,
    \qquad
    a=1,\ldots,k.
\end{equation}
Majoranas belonging to different tensor factors commute.  It will be
convenient to replace them by $2k$ mutually anti-commuting sets.

For each copy define the total parity operator
\begin{equation}
    Q_{L/R}^{(a)}
    =
    (-1)^{
       \sum_{j=1}^{N/2}N_j^{(a)}
    }
    =
    \prod_{j=1}^{N/2}
    Z_{j,L/R}^{(a)}
    =
    (-i)^{N/2}
    \prod_{l=1}^N
    \chi_{l,L/R}^{(a)}.
    \label{eq:SYK_copy_parity}
\end{equation}
For even $N$,
\begin{equation}
    \{
       Q_{L/R}^{(a)},
       \chi_{l,L/R}^{(a)}
    \}
    =
    0,
    \qquad
    (Q_{L/R}^{(a)})^2=1,
\end{equation}
while $Q_{L/R}^{(a)}$ commutes with Majoranas on all other copies.

Define
\begin{equation}
    \Sigma^{(a)}
    :=
    \left(
       \prod_{b=1}^{a-1}
       Q_L^{(b)}Q_R^{(b)}
    \right)
    Q_L^{(a)},
\end{equation}
and introduce
\begin{equation}
    \psi_{l,L}^{(a)}
    :=
    i\Sigma^{(a)}
    \chi_{l,L}^{(a)},
    \qquad
    \psi_{l,R}^{(a)}
    :=
    \Sigma^{(a)}
    \chi_{l,R}^{(a)}.
    \label{eq:SYK_psi_def}
\end{equation}
One checks that these are self-adjoint Majorana fermions satisfying
\begin{equation}
    \{
       \psi_{l,\alpha}^{(a)},
       \psi_{l',\beta}^{(b)}
    \}
    =
    2\delta_{ll'}\delta_{ab}\delta_{\alpha\beta},
    \qquad
    \alpha,\beta=L,R.
    \label{eq:majorana_CAR}
\end{equation}

A useful feature of this definition is that for even $q$ the Klein
factors cancel inside every $q$-body interaction:
\begin{align}
    \chi_{i_1,L}^{(a)}
    \cdots
    \chi_{i_q,L}^{(a)}
    &=
    \psi_{i_1,L}^{(a)}
    \cdots
    \psi_{i_q,L}^{(a)},
    \nonumber\\
    \chi_{i_1,R}^{(a)}
    \cdots
    \chi_{i_q,R}^{(a)}
    &=
    \psi_{i_1,R}^{(a)}
    \cdots
    \psi_{i_q,R}^{(a)}.
    \label{eq:SYK_evenq_Klein}
\end{align}

Define
\begin{equation}
    \Psi_{I,\alpha}^{(a)}
    :=
    \psi_{i_1,\alpha}^{(a)}
    \cdots
    \psi_{i_q,\alpha}^{(a)}.
\end{equation}
Using Eq.~\eqref{eq:SYK_W_transpose}, the hatted interaction appearing
in the $k$-twirl is
\begin{align}
    \widehat W_I^{(k)}
    &=
    i^{q/2}
    \sum_{a=1}^k
    \Psi_{I,L}^{(a)}
    -
    (-i)^{q/2}
    \sum_{a=1}^k
    \Psi_{I,R}^{(a)}.
    \label{eq:VSYK}
\end{align}

It is often convenient to bundle the replica label
$a=1,\ldots,k$ and the side label $\alpha\in\{L,R\}$ into a single
flavor index
\begin{equation}
    A=(a,\alpha),
    \qquad
    A=1,\ldots,2k.
\end{equation}
Define
\begin{equation}
    s_A
    =
    \begin{cases}
       1, & \alpha=L,\\[1mm]
       -(-1)^{q/2}, & \alpha=R.
    \end{cases}
    \label{eq:SYK_sA}
\end{equation}
Then
\begin{equation}
    \widehat W_I^{(k)}
    =
    i^{q/2}
    \sum_{A=1}^{2k}
    s_A\Psi_I^A.
    \label{eq:SYK_hatW_flavor}
\end{equation}


\subsection{$k$-twirl map in large $N$ SYK}

From Sec.~\ref{subsec:quenched_syk_gaussian_regulator}, the effective
operator on the doubled moment space is
\begin{equation}
    \widehat G_{\rm eff}^{(k)}
    =
    \frac{\sigma}{2}
    \sum_{i_1<\cdots<i_q}
    \left(
       \widehat W_{i_1\cdots i_q}^{(k)}
    \right)^2.
    \label{eq:SYK_Geff_start}
\end{equation}

Motivated by the form of Eq.~\eqref{eq:SYK_hatW_flavor}, define the
antisymmetric bilinears
\begin{equation}
    M_i^{AB}
    :=
    \frac{i}{2}
    [\psi_i^A,\psi_i^B],
    \qquad
    M_i^{AB}=-M_i^{BA}.
    \label{eq:def_M}
\end{equation}
For $A\neq B$,
\begin{equation}
    M_i^{AB}
    =
    i\psi_i^A\psi_i^B,
    \qquad
    (M_i^{AB})^\dagger=M_i^{AB},
    \qquad
    (M_i^{AB})^2=1,
    \label{eq:M_as_product}
\end{equation}
while
\begin{equation}
    M_i^{AA}=0.
\end{equation}

For $i\neq j$, the bilinears commute:
\begin{equation}
    [M_i^{AB},M_j^{CD}]
    =
    0,
    \qquad
    i\neq j.
    \label{eq:M_commute_different_sites}
\end{equation}
For the same site,
\begin{equation}
    [M_i^{AB},\psi_j^C]
    =
    2i\delta_{ij}
    \left(
       \delta^{BC}\psi_i^A
       -
       \delta^{AC}\psi_i^B
    \right),
    \label{eq:M_comm_psi}
\end{equation}
and hence
\begin{equation}
    [M_i^{AB},M_j^{CD}]
    =
    2i\delta_{ij}
    \left(
       \delta^{BC}M_i^{AD}
       -
       \delta^{AC}M_i^{BD}
       -
       \delta^{BD}M_i^{AC}
       +
       \delta^{AD}M_i^{BC}
    \right).
    \label{eq:so2k_commutator}
\end{equation}
Thus each site carries a copy of $\mathfrak{so}(2k)$, while generators
at different sites commute.

For $A\neq B$ and even $q$,
\begin{equation}
    \Psi_I^A\Psi_I^B
    =
    M_{i_1}^{AB}\cdots M_{i_q}^{AB},
    \label{eq:SYK_Psi_product}
\end{equation}
while
\begin{equation}
    i^q(\Psi_I^A)^2=1.
\end{equation}
Substituting Eq.~\eqref{eq:SYK_hatW_flavor} into
Eq.~\eqref{eq:SYK_Geff_start} therefore gives
\begin{equation}
    \boxed{
    \widehat G_{\rm eff}^{(k)}
    =
    \sigma k\binom Nq
    +
    \sigma i^q
    \sum_{A<B}
    s_As_B\,e_q^{AB},
    }
    \label{eq:Geff_poly_form}
\end{equation}
where
\begin{equation}
    e_q^{AB}
    :=
    \sum_{i_1<\cdots<i_q}
    M_{i_1}^{AB}\cdots M_{i_q}^{AB}.
    \label{eq:eABq}
\end{equation}


\subsection{Properties of the generator $G_{\rm eff}^{(k)}$}

Fix a pair $A<B$ and define
\begin{equation}
    S^{AB}
    :=
    \sum_{i=1}^N M_i^{AB}.
    \label{eq:SYK_SAB}
\end{equation}
For fixed $(A,B)$, the operators $M_i^{AB}$ commute at distinct sites and
square to identity.  Therefore they are simultaneously diagonalizable
with eigenvalues $\pm1$, and
\begin{equation}
    S^{AB}
    \in
    \{-N,-N+2,\ldots,N-2,N\},
    \label{eq:SYK_S_spectrum}
\end{equation}
with binomial multiplicities.

Moreover, $e_q^{AB}$ is a polynomial in the single collective variable
$S^{AB}$.  Its generating function is
\begin{align}
    \sum_{q=0}^N z^qe_q^{AB}
    &=
    \prod_{i=1}^N
    (1+zM_i^{AB})
    \nonumber\\
    &=
    (1+z)^{(N+S^{AB})/2}
    (1-z)^{(N-S^{AB})/2}.
    \label{eq:SYK_e_generating}
\end{align}
For even $q$,
\begin{equation}
    e_q^{AB}
    =
    \sum_{r=0}^{q/2}
    c_{q,r}(N)
    (S^{AB})^{q-2r}.
    \label{eq:eABq_poly}
\end{equation}
For example,
\begin{equation}
    e_2^{AB}
    =
    \frac{
       (S^{AB})^2-N
    }{2},
\end{equation}
and
\begin{equation}
    e_4^{AB}
    =
    \frac{
       (S^{AB})^4
       -
       2(3N-4)(S^{AB})^2
       +
       3N(N-2)
    }{24}.
    \label{eq:SYK_e4}
\end{equation}

It is convenient to strip off factorials and define the monic
normalization of the Krawtchouk polynomial
\begin{equation}
    P_q(S)
    :=
    q!\,e_q(S).
\end{equation}
Newton identities give the three-term recursion
\begin{equation}
    P_{q+1}(S)
    =
    SP_q(S)
    -
    q(N-q+1)P_{q-1}(S),
    \qquad
    P_0(S)=1,
    \qquad
    P_1(S)=S.
    \label{eq:krawtchouk_recursion}
\end{equation}
Equivalently, $P_q(S)$ is the characteristic polynomial of the
$q\times q$ tridiagonal matrix with diagonal entries $S$ and off-diagonal
entries
\begin{equation}
    \sqrt{m(N-m+1)},
    \qquad
    m=1,\ldots,q-1.
\end{equation}

Expanding
\begin{equation}
    P_q(S)
    =
    \sum_{r=0}^{\lfloor q/2\rfloor}
    \widetilde c_{q,r}(N)S^{q-2r}
\end{equation}
gives
\begin{equation}
    c_{q,r}(N)
    =
    \frac{\widetilde c_{q,r}(N)}{q!},
\end{equation}
with recursion
\begin{equation}
    (q+1)c_{q+1,r}(N)
    =
    c_{q,r}(N)
    -
    (N-q+1)c_{q-1,r-1}(N).
    \label{eq:cqr_recursion}
\end{equation}
A fully explicit form is
\begin{equation}
    c_{q,r}(N)
    =
    \frac{(-1)^r}{q!}
    \sum_{\substack{
       1\leq m_1<\cdots<m_r\leq q-1\\
       m_{j+1}\geq m_j+2
    }}
    \prod_{j=1}^r
    m_j(N-m_j+1).
    \label{eq:cqr_matching}
\end{equation}
In particular,
\begin{equation}
    c_{q,0}(N)=\frac1{q!},
    \qquad
    c_{q,1}(N)
    =
    -
    \frac{
       q(q-1)(3N-2q+4)
    }{
       6q!
    },
\end{equation}
and for $q=2n$ and even $N$,
\begin{equation}
    c_{2n,n}(N)
    =
    (-1)^n
    \binom{N/2}{n}.
\end{equation}


\subsubsection{Large-$N$ limits of $e_q^{AB}$}

Define
\begin{equation}
    m^{AB}
    :=
    \frac{S^{AB}}{N}.
\end{equation}

\paragraph{1) $N\rightarrow\infty$ with $q$ fixed.}

On each joint eigenspace,
\begin{equation}
    \frac{e_q^{AB}}{\binom Nq}
    =
    (m^{AB})^q
    -
    \frac{q(q-1)}{2N}
    \left(
       1-(m^{AB})^2
    \right)
    (m^{AB})^{q-2}
    +
    O(N^{-2}).
    \label{eq:eq_largeN_fixedq}
\end{equation}
In particular,
\begin{equation}
    e_q^{AB}
    =
    \binom Nq
    (m^{AB})^q
    +
    O(N^{q-1}).
\end{equation}

\paragraph{2) $N,q\rightarrow\infty$ with $q^2/N$ fixed.}

Let
\begin{equation}
    \frac{q^2}{N}
    =
    \lambda
\end{equation}
be held fixed.  The generating function
Eq.~\eqref{eq:SYK_e_generating} gives
\begin{align}
    &(1+z)^{(N+S^{AB})/2}
     (1-z)^{(N-S^{AB})/2}
    \nonumber\\
    &\hspace{15mm}
    =
    \exp\left[
       S^{AB}z
       -
       \frac N2z^2
       +
       O(S^{AB}z^3)
       +
       O(Nz^4)
    \right].
\end{align}
For the coefficient regime $q=O(\sqrt N)$, the leading quadratic
approximation gives the Hermite scaling
\begin{equation}
    e_q^{AB}
    \sim
    \frac{N^{q/2}}{q!}
    {\rm He}_q
    \left(
       \frac{S^{AB}}{\sqrt N}
    \right),
    \label{eq:eq_hermite_scaling}
\end{equation}
where the probabilists' Hermite polynomials are defined by
\begin{equation}
    e^{xt-t^2/2}
    =
    \sum_{q=0}^\infty
    {\rm He}_q(x)
    \frac{t^q}{q!}.
\end{equation}

Using Eq.~\eqref{eq:eABq_poly}, the doubled-space generator may therefore
also be written
\begin{equation}
    \widehat G_{\rm eff}^{(k)}
    =
    \sigma k\binom Nq
    +
    \sigma i^q
    \sum_{A<B}
    s_As_B
    \sum_{r=0}^{q/2}
    c_{q,r}(N)
    (S^{AB})^{q-2r}.
    \label{eq:Geff_sum_over_AB}
\end{equation}
At each site $i$, the $M_i^{AB}$ generate a copy of
$\mathfrak{so}(2k)$, while different sites commute.  Therefore the global
sums $S^{AB}=\sum_iM_i^{AB}$ generate the diagonal
$\mathfrak{so}(2k)$ action, and
Eq.~\eqref{eq:Geff_sum_over_AB} exhibits
$\widehat G_{\rm eff}^{(k)}$ as an element of
$U(\mathfrak{so}(2k))$.

Finally, on the physical $k$-replica Hilbert space define the Hermitian
operators
\begin{equation}
    X_I
    :=
    \sqrt{\sigma}
    \sum_{a=1}^kW_I^{(a)}.
\end{equation}
The corresponding positive generator is
\begin{align}
    G_{\rm eff}^{(k)}(Y)
    &=
    \frac12
    \sum_I
    [X_I,[X_I,Y]]
    \nonumber\\
    &=
    -\sum_I
    \left(
       X_IYX_I
       -
       \frac12\{X_I^2,Y\}
    \right).
    \label{eq:SYK_superoperator_G}
\end{align}
Thus $-G_{\rm eff}^{(k)}$ is in Lindblad form with self-adjoint jump
operators $X_I$, and the large-$N$ quenched $k$-twirl is
\begin{equation}
    \boxed{
    {\cal E}_t^{(k)}
    \simeq
    e^{-t^2G_{\rm eff}^{(k)}}.
    }
\end{equation}


    \section{Exponential decay and semi-groups from Brownian dynamics}

    To simplify our notation, from now on, we use dimensionless parameters and denote
    them by $\tilde{t}\to t$, $\tilde{\sigma}\to \sigma$, $\tilde{g}_{I}\to g_{I}$,
    and $\tilde{H}\to H$. We repeat the example above, but this time, splitting the
    total evolution time $T$ into $T/\epsilon$ intervals $(t_{i},t_{i}+\epsilon)$,
    and making the coupling to depend on the interval $g(t_{i})$:
    \begin{eqnarray}
        H_g(t_i)=\sum_I g_I(t_i)V_I\ .
    \end{eqnarray}
    We pick the coupling $g(t_{i})$ during this time interval by taking i.i.d.
    samples from a Gaussian distribution so that
    \begin{eqnarray}
        &&\mathcal{E}_g(g_I(t_i))=0\nn\\
        &&\mathcal{E}_g(g_I(t_i)g_J(t_j))=\frac{1}{\ep}\delta_{ij}\sigma
    \end{eqnarray}
    After one time interval $\ep$, we have have
    \begin{eqnarray}
        &&\mathcal{E}_g(dU^{\otimes k}\otimes (dU^*)^{\otimes k})=\mathcal{E}_g\lb
        1^{\otimes 2k}-i\ep \hat{\mathbb{H}}_g-\frac{\ep^{2}}{2}\hat{\mathbb{H}}_g^2+\cdots\rb
        \nn\\
        &&dU=U(t_i,t_i+\ep),\qquad \hat{\mathbb{H}}_g=\sum_{a=1}^k \hat{H}_g^{(a)},\qquad\hat{H}_g^{(a)}=H_{g,L}^{(a)}-(H_{g,R}^{(a)})^T\nn\ .
    \end{eqnarray}
    where the index $a=1,\cdots , k$ keeps track of the $k$ copies in the twirl map.
    Note that since different samples of $g_{I}(t_{i})$ result in operators that
    do not commute, we cannot pick a vector $\ket{1}$ that simulatenously diagonalizes
    them all. Therefore, for any fixed basis, we need to write
    $(H_{g,R}^{(a)})^{T}$ as the mirror operator of $H_{g,L}^{(a)}$.

    Since
    \begin{eqnarray}
        &&\mathcal{E}_g(\hat{\mathbb{H}_g})=\sum_{a I}\mathcal{E}(g_I)\hat{V}_I^{(a)}=0\nn\\
        &&\mathcal{E}_g(\hat{\mathbb{H}}_g^2)=\sum_{a,a',I,J}\mathcal{E}_g(g_Ig_J)\hat{V}^{(a)}_I
        \hat{V}_J^{(a')}=\frac{\sigma}{\ep}\sum_{a,a',I} \hat{V}_I^{(a)}\hat{V}_I^{(a')}\ .
    \end{eqnarray}
    We define
    \begin{eqnarray}
        &&\hat{\mathbb{G}}_{eff}^{(k)}\equiv \ep\: \mathcal{E}_g\lb \hat{\mathbb{H}}_g^2\rb=\sum_I
        \hat{\mathbb{X}}_I^2\nn\\
        &&\hat{\mathbb{X}}_I=\sqrt{\sigma}\sum_a \hat{V}_I^{(a)}=\mathbb{X}_{I,L}-\mathbb{X}_{I,R}^T\nn\\
        &&\mathbb{X}_{I}=\sqrt{\sigma}\sum_a V_I^{(a)}\ .
    \end{eqnarray}
    It follows that
    \begin{eqnarray}
        \mathcal{E}_g(dU^{\otimes k}\otimes (dU^*)^{\otimes k})= 1^{\otimes 2k}-\frac{\ep}{2}\hat{\mathbb{G}}_{eff}^{(k)}+O(\ep^2)\ .
    \end{eqnarray}
    Then, we find
    \begin{eqnarray}
        &&\mathcal{E}_g(dU^{\otimes k}\otimes (dU^*)^{\otimes k})(\mathbb{Y}\otimes
        1^{\otimes k})\ket{1}^{\otimes k}\nn\\
        &&=\lb 1^{\otimes 2k}-\frac{\ep}{2}\hat{\mathbb{G}}_{eff}^{(k)}\rb(\mathbb{Y}\otimes
        1^{\otimes k})\ket{1}^{\otimes k}\nn\\
        &&=\lb \lb 1-\ep \:\mathcal{G}_{eff}^{(k)}\rb(\mathbb{Y})\otimes 1^{\otimes 1}\rb\ket{1}^{\otimes k}
    \end{eqnarray}
    where $\mathbb{Y}$ is a general operator acting on $\mathcal{H}^{\otimes k}$,
    and the super-operator $\mathcal{G}$ generates a Lindblad evolution
    \begin{eqnarray}
        \mathcal{E}_g(\p_t \mathbb{Y})=-\mathcal{G}_{eff}^{(k)}(\mathbb{Y})&=&\sum_{I,a,a'}\lb
        V_I^{(a)} \mathbb{Y} V_I^{(a')}-\frac{1}{2}\{ V_I^{(a)}V_I^{(a')},\mathbb{Y}\}\rb\nn\\
        &=&\sum_I \lb \mathbb{X}_I Y \mathbb{X}_I-\frac{1}{2}\{\mathbb{X}_I^2,\mathbb{Y}\}\rb\nn\\
        &=&-\frac{1}{2}\sum_I [\mathbb{X}_I,[\mathbb{X}_I,\mathbb{Y}]]\ .
    \end{eqnarray}
    where we have used $[V_{I}^{(a)},V_{I}^{(a')}]=0$. Here, our Lindblad jump
    operators are the self-adjoint $\mathbb{X}_{I}$ Lindblad operators. Alternatively,
    we can write this Lindbladian in the Schrodinger picture for the density
    matrix by replacing $Y$ with $\rho^{\otimes k}$:
    \begin{eqnarray}
        \mathcal{E}_g(\p_t \rho^{\otimes k})=-\mathcal{G}_{eff}^{(k)}(\rho^{\otimes k})\ .
    \end{eqnarray}

    Integrating our Lindbladian evolution over a time-interval $(t,t+T)$ results
    in a Markov semi-group
    \begin{eqnarray}
        \mathcal{E}_{g;T}(\mathbb{Y})=e^{-T\mathcal{G}_{eff}^{(k)}}(\mathbb{Y})\ .
    \end{eqnarray}

 

    \section{Semigroups in Brownian bosonic and Spin SYK models}
\label{sec:brownian_spin_syk}

In Sec.~6 we considered a general time-dependent Hamiltonian
\begin{equation}
    H_g(t)=\sum_I g_I(t)V_I
\end{equation}
whose couplings are independent Gaussian white-noise processes,
\begin{equation}
    \mathbb E_g[g_I(t)]=0,
    \qquad
    \mathbb E_g[g_I(t)g_J(t')]
    =
    \sigma\,\delta_{IJ}\delta(t-t').
    \label{eq:spinbrown_white_noise}
\end{equation}
Splitting the evolution into intervals of width $\epsilon$, the variance of
$g_I$ on one interval scales as $\sigma/\epsilon$.  Consequently only the
quadratic term in the infinitesimal evolution survives as
$\epsilon\rightarrow0$, while the higher moments contribute at higher powers
of $\epsilon$.  The disorder-averaged $k$-twirl is therefore an exact Markov
semigroup.

For later reference, let
\begin{equation}
    W_I^{(k)}
    :=
    \sum_{a=1}^kV_I^{(a)}
\end{equation}
on the physical $k$-replica Hilbert space and
\begin{equation}
    \widehat W_I^{(k)}
    :=
    \sum_{a=1}^k
    \left(
        V_{I,L}^{(a)}
        -
        (V_{I,R}^{(a)})^T
    \right)
\end{equation}
on the doubled space.  The result of Sec.~6 can be written as
\begin{equation}
    \widehat G_{\rm eff}^{(k)}
    =
    \frac{\sigma}{2}
    \sum_I
    \left(
        \widehat W_I^{(k)}
    \right)^2,
    \label{eq:spinbrown_Ghat_general}
\end{equation}
or, equivalently, as the superoperator
\begin{align}
    G_{\rm eff}^{(k)}(Y)
    &=
    \frac{\sigma}{2}
    \sum_I
    [W_I^{(k)},[W_I^{(k)},Y]]
    \nonumber\\
    &=
    -\sigma
    \sum_I
    \left(
       W_I^{(k)}YW_I^{(k)}
       -
       \frac12
       \{
          (W_I^{(k)})^2,Y
       \}
    \right).
    \label{eq:spinbrown_G_general}
\end{align}
Thus
\begin{equation}
    \boxed{
    {\cal E}_{g;T}^{(k)}
    =
    e^{-T G_{\rm eff}^{(k)}} .
    }
    \label{eq:spinbrown_semigroup_general}
\end{equation}

Notice the important distinction from the quenched calculation of Sec.~4.
For a time-independent Gaussian Hamiltonian the analogous quadratic
generator emerged only after taking the large $N$ limit and suppressing
non-commuting terms in the ordered Wick expansion.  Here
Eq.~\eqref{eq:spinbrown_semigroup_general} is exact at finite $N$.  The
difference also appears in the time dependence: quenched Gaussian averaging
produces an exponent proportional to $t^2$, whereas white-noise averaging
produces a semigroup whose exponent is proportional to $T$.


\subsection{Brownian three-component Spin SYK}
\label{subsec:brownian_swingle}

We first apply Eq.~\eqref{eq:spinbrown_semigroup_general} to the bosonic
Spin SYK model of \cite{SwingleWiner2024}.  There are $N$ spin-$1/2$ sites,
and the interaction operators are
\begin{equation}
    V_{I}
    =
    \sigma_{r_1}^{\mu_1}
    \cdots
    \sigma_{r_q}^{\mu_q},
    \qquad
    r_1<\cdots<r_q,
    \qquad
    \mu_j\in\{x,y,z\}.
    \label{eq:spinbrown_SW_VI}
\end{equation}
There are
\begin{equation}
    M_q^{(3)}
    =
    3^q\binom{N}{q}
\end{equation}
such operators.  We Brownianize the model by taking
\begin{equation}
    H_{\rm SW}(t)
    =
    \sum_{\substack{r_1<\cdots<r_q\\
                    \mu_1,\ldots,\mu_q}}
    J_{r_1\mu_1\cdots r_q\mu_q}(t)
    \sigma_{r_1}^{\mu_1}\cdots\sigma_{r_q}^{\mu_q},
    \label{eq:spinbrown_SW_H}
\end{equation}
with
\begin{equation}
    \mathbb E[J_I(t)J_J(t')]
    =
    \sigma_q^{(3)}
    \delta_{IJ}\delta(t-t').
    \label{eq:spinbrown_SW_noise}
\end{equation}
If we preserve the $N$-scaling of the static Spin SYK normalization, it is
natural to write
\begin{equation}
    \sigma_q^{(3)}
    =
    \frac{(q-1)!{\cal J}_3}{N^{q-1}} .
    \label{eq:spinbrown_SW_sigma}
\end{equation}
The Brownian parameter ${\cal J}_3$ has the dimensions appropriate to a
white-noise variance density and should not be identified with the static
coupling without first choosing units of time.

Defining
\begin{equation}
    W_{r_1\mu_1\cdots r_q\mu_q}
    :=
    \sum_{a=1}^k
    \left(
      \sigma_{r_1}^{\mu_1}\cdots
      \sigma_{r_q}^{\mu_q}
    \right)^{(a)},
\end{equation}
the exact generator is
\begin{equation}
    \boxed{
    G_{{\rm SW},q}^{(k)}(Y)
    =
    \frac{\sigma_q^{(3)}}{2}
    \sum_{\substack{r_1<\cdots<r_q\\
                    \mu_1,\ldots,\mu_q}}
    [W_{r_1\mu_1\cdots r_q\mu_q},
     [W_{r_1\mu_1\cdots r_q\mu_q},Y]] .
    }
    \label{eq:spinbrown_SW_Gq}
\end{equation}
Hence
\begin{equation}
    {\cal E}_{{\rm SW};T}^{(k)}
    =
    e^{-T G_{{\rm SW},q}^{(k)}} .
    \label{eq:spinbrown_SW_semigroup}
\end{equation}

For $q=2$,
\begin{equation}
    G_{{\rm SW},2}^{(k)}(Y)
    =
    \frac{{\cal J}_3}{2N}
    \sum_{r<s}
    \sum_{\mu,\nu=x,y,z}
    [W_{rs}^{\mu\nu},
     [W_{rs}^{\mu\nu},Y]],
    \label{eq:spinbrown_SW_q2}
\end{equation}
where
\begin{equation}
    W_{rs}^{\mu\nu}
    =
    \sum_{a=1}^k
    (\sigma_r^\mu\sigma_s^\nu)^{(a)}.
\end{equation}
For $q=4$,
\begin{equation}
    G_{{\rm SW},4}^{(k)}(Y)
    =
    \frac{3{\cal J}_3}{N^3}
    \sum_{r_1<r_2<r_3<r_4}
    \sum_{\mu_1,\ldots,\mu_4=x,y,z}
    [W_I,[W_I,Y]] .
    \label{eq:spinbrown_SW_q4}
\end{equation}


\subsubsection{The four-replica sector and operator-size dynamics}

For the second moment, $k=2$, the doubled Hilbert space contains four
replicas.  After reordering them as forward, backward, forward, backward,
Eq.~\eqref{eq:spinbrown_Ghat_general} becomes
\begin{equation}
    {\cal H}_{4}
    =
    \frac{\sigma_q^{(3)}}{2}
    \sum_A
    \left(
       V_A^{a}
       -
       V_A^{*,b}
       +
       V_A^{c}
       -
       V_A^{*,d}
    \right)^2.
    \label{eq:spinbrown_SW_four_replica}
\end{equation}
This is the four-replica Brownian Hamiltonian used in
\cite{Xu2024}.  Thus the Brownian version of the Swingle-Winer model is
precisely the Brownian spin model studied there.

Let $S$ be a Pauli string.  Since $V_A^2=1$, the contribution of a single
interaction to the diagonal four-replica sector is zero when
$[V_A,S]=0$, while, when $\{V_A,S\}=0$,
\begin{equation}
    \frac12
    \left(
       V_A^{a}
       -
       V_A^{*,b}
       +
       V_A^{c}
       -
       V_A^{*,d}
    \right)^2
    |S\otimes S^\dagger\rangle
    =
    4|S\otimes S^\dagger\rangle
    -
    4|V_AS\otimes S^\dagger V_A\rangle .
    \label{eq:spinbrown_transition_action}
\end{equation}
Therefore each interaction which anticommutes with $S$ generates a
classical transition
\begin{equation}
    S\longrightarrow V_AS
\end{equation}
at rate
\begin{equation}
    4\sigma_q^{(3)}.
    \label{eq:spinbrown_single_transition_rate}
\end{equation}

Because the Swingle-Winer ensemble contains \emph{all} Pauli strings of
size $q$, all Pauli strings $S$ of the same size are equivalent under the
ensemble.  Consequently the stochastic evolution closes on the ordinary
operator size
\begin{equation}
    m=|{\rm supp}(S)|.
\end{equation}
Without loss of generality take
\begin{equation}
    S=\sigma_1^z\cdots\sigma_m^z.
\end{equation}
Parameterize the size change as
\begin{equation}
    \Delta m=q-2\ell+1,
    \qquad
    \ell=1,\ldots,q.
\end{equation}
The exact transition rate is
\begin{align}
    \gamma_\ell^{(3)}(m)
    =
    4\sigma_q^{(3)}
    \sum_{n_z}
    &
    2^{\,2\ell-2n_z-1}
    3^{\,n_z+q-2\ell+1}
    \binom{m}{n_z}
    \binom{N-m}{n_z+q-2\ell+1}
    \nonumber\\
    &\times
    \binom{m-n_z}{2\ell-2n_z-1},
    \label{eq:spinbrown_SW_gamma_general}
\end{align}
where a binomial coefficient with an inadmissible lower argument is
understood to vanish.  The exact master equation is therefore
\begin{equation}
    \boxed{
    \partial_T P(m,T)
    =
    \sum_{\ell=1}^q
    \left[
       \gamma_\ell^{(3)}
       (m+2\ell-1-q)
       P(m+2\ell-1-q,T)
       -
       \gamma_\ell^{(3)}(m)P(m,T)
    \right].
    }
    \label{eq:spinbrown_SW_master}
\end{equation}
This is exactly the spin master equation derived in \cite{Xu2024}.


\paragraph{Quadratic model.}

For $q=2$, there are only two possible transitions,
\begin{equation}
    \Delta m=+1,-1,
\end{equation}
with rates
\begin{align}
    \gamma_1^{(3)}(m)
    &=
    24\sigma_2^{(3)}m(N-m),
    \label{eq:spinbrown_SW_q2_up}\\
    \gamma_2^{(3)}(m)
    &=
    8\sigma_2^{(3)}m(m-1).
    \label{eq:spinbrown_SW_q2_down}
\end{align}
Thus
\begin{align}
    \partial_TP(m)
    =&\,
    \gamma_1^{(3)}(m-1)P(m-1)
    +
    \gamma_2^{(3)}(m+1)P(m+1)
    \nonumber\\
    &-
    \left[
       \gamma_1^{(3)}(m)
       +
       \gamma_2^{(3)}(m)
    \right]P(m).
    \label{eq:spinbrown_SW_q2_master}
\end{align}


\paragraph{Quartic model.}

For $q=4$,
\begin{equation}
    \Delta m=3,1,-1,-3.
\end{equation}
The four rates following from
Eq.~\eqref{eq:spinbrown_SW_gamma_general} are
\begin{align}
    \gamma_1^{(3)}(m)
    &=
    216\sigma_4^{(3)}
    m\binom{N-m}{3},
    \label{eq:spinbrown_SW_q4_g1}\\
    \gamma_2^{(3)}(m)
    &=
    96\sigma_4^{(3)}
    \binom{m}{3}(N-m)
    +
    72\sigma_4^{(3)}
    m(m-1)\binom{N-m}{2},
    \label{eq:spinbrown_SW_q4_g2}\\
    \gamma_3^{(3)}(m)
    &=
    32\sigma_4^{(3)}
    m\binom{m-1}{3}
    +
    24\sigma_4^{(3)}
    \binom{m}{2}(m-2)(N-m),
    \label{eq:spinbrown_SW_q4_g3}\\
    \gamma_4^{(3)}(m)
    &=
    32\sigma_4^{(3)}
    \binom{m}{4}.
    \label{eq:spinbrown_SW_q4_g4}
\end{align}


\subsubsection{Comparison with the large $N$ solution of Xu}

The Brownian clock used in \cite{Xu2024} is fixed by choosing
\begin{equation}
    \sigma_{q,{\rm Xu}}^{(3)}
    =
    \frac{N}{
       8q\,3^{q-1}\binom{N}{q}
    }.
    \label{eq:spinbrown_Xu_normalization}
\end{equation}
With this normalization the largest-growth transition has the simple
large-$N$ limit
\begin{equation}
    \gamma_1^{(3)}(m)\longrightarrow m,
    \qquad
    N\rightarrow\infty
\end{equation}
for fixed $m$, while all other transition rates are suppressed.  Hence
Eq.~\eqref{eq:spinbrown_SW_master} reduces to
\begin{equation}
    \partial_TP(m)
    =
    (m-q+1)P(m-q+1)-mP(m).
    \label{eq:spinbrown_SW_early_master}
\end{equation}
The elementary growth step is $q-1$, and the early-time growth exponent is
\begin{equation}
    \lambda_{\rm SW}=q-1.
\end{equation}
Thus
\begin{equation}
    \lambda_{\rm SW}
    =
    \begin{cases}
    1, & q=2,\\
    3, & q=4.
    \end{cases}
\end{equation}

More generally, \cite{Xu2024} solves the large-$N$ size distribution at all
times.  In the normalization
Eq.~\eqref{eq:spinbrown_Xu_normalization}, define
\begin{equation}
    x=\frac{m}{N},
    \qquad
    x_{\rm eq}=\frac34,
    \qquad
    \lambda=q-1,
\end{equation}
and
\begin{equation}
    T_*
    =
    \frac{1}{\lambda}
    \left(
       \log N-\log\lambda
    \right).
\end{equation}
For an initial operator of size $m_0$, the large-$N$ result takes the form
\begin{equation}
    p(x,T)
    =
    \frac{dz}{dx}
    \frac{
       e^{-z}z^{m_0/\lambda-1}
    }{
       \Gamma(m_0/\lambda)
    },
    \label{eq:spinbrown_Xu_distribution}
\end{equation}
where
\begin{equation}
    z(x,T)
    =
    \frac{x_{\rm eq}}{\lambda}
    e^{\lambda(T_*-T)}
    \left[
       -1+
       \left(
          1-\frac{x}{x_{\rm eq}}
       \right)^{-\lambda}
    \right].
    \label{eq:spinbrown_Xu_z}
\end{equation}
Therefore all of the large-$N$ results of \cite{Xu2024} apply directly to
the Brownianized Swingle-Winer model.  If instead we use the normalization
Eq.~\eqref{eq:spinbrown_SW_sigma}, the only change is an overall rescaling
of the Brownian clock,
\begin{equation}
    T_{\rm Xu}
    =
    \frac{
       \sigma_q^{(3)}
    }{
       \sigma_{q,{\rm Xu}}^{(3)}
    }T .
    \label{eq:spinbrown_time_rescale}
\end{equation}


\subsection{Brownian two-component Spin SYK}
\label{subsec:brownian_BDN}

We next consider the Spin SYK model of \cite{BasuDasNandy2026}.  Denote the
number of physical spin sites by
\begin{equation}
    n:=N_{\rm Spin}.
\end{equation}
The $2n$ constituent operators are
\begin{equation}
    O_{2r-1}=\sigma_r^x,
    \qquad
    O_{2r}=\sigma_r^y,
    \qquad
    r=1,\ldots,n.
    \label{eq:spinbrown_BDN_O}
\end{equation}
The full Brownian Hamiltonian is
\begin{equation}
    H_{{\rm Spin},q}(t)
    =
    \sum_{1\leq i_1<\cdots<i_q\leq2n}
    J_{i_1\cdots i_q}(t)
    V_{i_1\cdots i_q},
    \label{eq:spinbrown_BDN_full_H}
\end{equation}
where
\begin{equation}
    V_{i_1\cdots i_q}
    =
    i^{\eta_{i_1\cdots i_q}}
    O_{i_1}\cdots O_{i_q}
    =
    V_{i_1\cdots i_q}^\dagger
\end{equation}
and
\begin{equation}
    \mathbb E[
       J_I(t)J_J(t')
    ]
    =
    \sigma_q^{(2)}
    \delta_{IJ}\delta(t-t').
\end{equation}
Keeping the normalization inherited from the static model, up to an
arbitrary choice of Brownian clock, we write
\begin{equation}
    \sigma_q^{(2)}
    =
    \frac{(q-1)!{\cal J}_2}
         {(2n)^{q-1}}.
    \label{eq:spinbrown_BDN_sigma}
\end{equation}

The exact $k$-twirl generator of the full model is
\begin{equation}
    \boxed{
    G_{{\rm full},q}^{(k)}(Y)
    =
    \frac{\sigma_q^{(2)}}{2}
    \sum_{1\leq i_1<\cdots<i_q\leq2n}
    [W_I,[W_I,Y]],
    }
    \label{eq:spinbrown_BDN_full_G}
\end{equation}
with
\begin{equation}
    W_I
    =
    \sum_{a=1}^kV_I^{(a)}.
\end{equation}

The genuine model, gSpin-SYK, removes all terms containing two constituent
operators on the same physical site.  Its interactions are simply
\begin{equation}
    V_{S,\boldsymbol\alpha}
    =
    \prod_{r\in S}\sigma_r^{\alpha_r},
    \qquad
    |S|=q,
    \qquad
    \alpha_r\in\{x,y\},
    \label{eq:spinbrown_BDN_genuine_V}
\end{equation}
and
\begin{equation}
    \boxed{
    G_{g,q}^{(k)}(Y)
    =
    \frac{\sigma_q^{(2)}}{2}
    \sum_{\substack{|S|=q\\
                    \boldsymbol\alpha\in\{x,y\}^q}}
    [W_{S,\boldsymbol\alpha},
     [W_{S,\boldsymbol\alpha},Y]] .
    }
    \label{eq:spinbrown_BDN_genuine_G}
\end{equation}
In both cases the evolution is the exact semigroup
\begin{equation}
    {\cal E}_{T}^{(k)}
    =
    e^{-T G^{(k)}}.
\end{equation}


\subsubsection{Quadratic model}

For the genuine quadratic model,
\begin{equation}
    G_{g,2}^{(k)}(Y)
    =
    \frac{{\cal J}_2}{4n}
    \sum_{r<s}
    \sum_{\alpha,\beta=x,y}
    [W_{rs}^{\alpha\beta},
     [W_{rs}^{\alpha\beta},Y]].
    \label{eq:spinbrown_BDN_q2_genuine}
\end{equation}
The full quadratic Hamiltonian additionally contains the $n$ same-site
terms
\begin{equation}
    iO_{2r-1}O_{2r}
    =
    \pm\sigma_r^z.
\end{equation}
Since the sign drops out of a double commutator,
\begin{equation}
    G_{{\rm full},2}^{(k)}
    =
    G_{g,2}^{(k)}
    +
    \frac{{\cal J}_2}{4n}
    \sum_{r=1}^n
    [W_r^z,[W_r^z,\cdot]],
    \label{eq:spinbrown_BDN_q2_full}
\end{equation}
where
\begin{equation}
    W_r^z=\sum_a(\sigma_r^z)^{(a)}.
\end{equation}


\subsubsection{Quartic model}

For the genuine quartic model,
\begin{equation}
    G_{g,4}^{(k)}(Y)
    =
    \frac{3{\cal J}_2}{8n^3}
    \sum_{r_1<r_2<r_3<r_4}
    \sum_{\alpha_1,\ldots,\alpha_4=x,y}
    [W_I,[W_I,Y]] .
    \label{eq:spinbrown_BDN_q4_genuine}
\end{equation}
The full model contains three classes of terms.  According to the number
$r$ of physical sites on which both $O_{2j-1}$ and $O_{2j}$ occur,
\begin{equation}
    r=0,1,2.
\end{equation}
After simplifying the Pauli products these are
\begin{equation}
    \begin{array}{c|c|c}
    r & \text{operator type} & \text{number of terms}\\ \hline
    0 &
    \sigma^{\alpha_1}\sigma^{\alpha_2}
    \sigma^{\alpha_3}\sigma^{\alpha_4}
    &
    16\binom n4
    \\[2mm]
    1 &
    \sigma^z\sigma^\alpha\sigma^\beta
    &
    4n\binom{n-1}{2}
    \\[2mm]
    2 &
    \sigma^z\sigma^z
    &
    \binom n2 .
    \end{array}
    \label{eq:spinbrown_BDN_q4_classes}
\end{equation}
Indeed,
\begin{equation}
    16\binom n4
    +
    4n\binom{n-1}{2}
    +
    \binom n2
    =
    \binom{2n}{4}.
\end{equation}
Consequently,
\begin{align}
    G_{{\rm full},4}^{(k)}
    =
    G_{g,4}^{(k)}
    +
    \frac{3{\cal J}_2}{8n^3}
    \bigg[
       &
       \sum_r
       \sum_{\substack{s<t\\s,t\neq r}}
       \sum_{\alpha,\beta=x,y}
       [W_{r;st}^{z;\alpha\beta},
        [W_{r;st}^{z;\alpha\beta},\cdot]]
       \nonumber\\
       &+
       \sum_{r<s}
       [W_{rs}^{zz},
        [W_{rs}^{zz},\cdot]]
    \bigg].
    \label{eq:spinbrown_BDN_q4_full}
\end{align}
The first correction contains $O(n^3)$ terms and the second $O(n^2)$
terms, compared with $O(n^4)$ genuine quartic terms.  Therefore the
collision sectors are subleading in the fixed-$q$ large-$n$ dynamics.
This is the Brownian analogue of the counting found for the quenched
model.


\subsection{Operator-size dynamics in the two-component model}
\label{subsec:BDN_size_dynamics}

There is an important difference between
Eq.~\eqref{eq:spinbrown_BDN_genuine_G} and the three-component model.
Although the diagonal Pauli-string subspace is still invariant, Pauli
strings of the same ordinary size are no longer equivalent.

For example, consider the genuine $q=2$ model.  Starting from a single
$\sigma^z$, there are
\begin{equation}
    4(n-1)
\end{equation}
two-body $x/y$ interactions which anticommute with it and increase its
support.  Starting instead from a single $\sigma^x$, there are only
\begin{equation}
    2(n-1)
\end{equation}
such interactions.  Thus two strings of the same size $m=1$ have different
transition rates.  The one-dimensional operator-size basis used for the
three-component model is therefore not closed.

A minimal closed description keeps track of
\begin{equation}
    u
    =
    \#\{\text{$x$ or $y$ sites}\},
    \qquad
    v
    =
    \#\{\text{$z$ sites}\},
    \qquad
    m=u+v.
    \label{eq:spinbrown_BDN_uv}
\end{equation}
Consider a Pauli string with $(u,v)$ and a genuine $q$-body interaction.
Among the $q$ selected sites, let
\begin{align}
    a&=\text{number of identity sites},\nonumber\\
    b&=\text{number of $z$ sites},\nonumber\\
    c&=\text{number of transverse sites on which $V$ has the same
       $x/y$ component},\nonumber\\
    d&=\text{number of transverse sites on which $V$ has the opposite
       $x/y$ component}.
\end{align}
Thus
\begin{equation}
    a+b+c+d=q.
\end{equation}
A $z$ site anticommutes with either $x$ or $y$, while a transverse site
anticommutes only when the other transverse component is chosen.  Hence
\begin{equation}
    \{V,S\}=0
    \qquad\Longleftrightarrow\qquad
    b+d\ \text{is odd}.
    \label{eq:spinbrown_BDN_anticommute}
\end{equation}

The number of interactions of type $(a,b,c,d)$ is
\begin{equation}
    {\cal N}_{abcd}(u,v)
    =
    2^{a+b}
    \binom{n-u-v}{a}
    \binom{v}{b}
    \frac{u!}{
       c!\,d!\,(u-c-d)!
    }.
    \label{eq:spinbrown_BDN_count}
\end{equation}
The factor $2^{a+b}$ comes from the independent $x/y$ choice on identity
and $z$ sites.  The $x/y$ choice on a transverse site is fixed once it is
specified to be the same or opposite component.

Multiplication by the interaction sends
\begin{equation}
    (u,v)
    \longrightarrow
    {\cal T}_{abcd}(u,v)
    :=
    \left(
       u+a+b-c-d,\,
       v-b+d
    \right),
    \label{eq:spinbrown_BDN_transition}
\end{equation}
and the corresponding rate is
\begin{equation}
    \boxed{
    \Gamma_{abcd}(u,v)
    =
    4\sigma_q^{(2)}
    {\cal N}_{abcd}(u,v),
    \qquad
    b+d\ {\rm odd}.
    }
    \label{eq:spinbrown_BDN_rate}
\end{equation}
Consequently the exact stochastic equation for the genuine model is
\begin{align}
    \partial_TP(u,v;T)
    =
    \sum_{\substack{a+b+c+d=q\\b+d\ {\rm odd}}}
    \bigg[
       &
       \Gamma_{abcd}
       \left(
          {\cal T}_{abcd}^{-1}(u,v)
       \right)
       P\left(
          {\cal T}_{abcd}^{-1}(u,v);T
       \right)
       \nonumber\\
       &-
       \Gamma_{abcd}(u,v)P(u,v;T)
    \bigg],
    \label{eq:spinbrown_BDN_master_uv}
\end{align}
where terms whose inverse image contains an inadmissible occupation number
are defined to vanish.

The ordinary size changes by
\begin{equation}
    \Delta m
    =
    a-c.
    \label{eq:spinbrown_BDN_delta_m}
\end{equation}
The rate, however, depends separately on $u$ and $v$, which is why summing
Eq.~\eqref{eq:spinbrown_BDN_master_uv} over fixed $m=u+v$ does not produce
a closed equation for $P(m)$.


\subsubsection{Exact quadratic master equation}

For $q=2$, Eq.~\eqref{eq:spinbrown_BDN_master_uv} contains four transitions.
Writing
\begin{equation}
    n_0=n-u-v,
\end{equation}
they are
\begin{align}
    (u,v)&\longrightarrow(u,v+1),
    &
    \Gamma_{+}^{(\perp)}
    &=
    8\sigma_2^{(2)}u n_0,
    \label{eq:spinbrown_BDN_q2_rate1}\\
    (u,v)&\longrightarrow(u+2,v-1),
    &
    \Gamma_{+}^{(z)}
    &=
    16\sigma_2^{(2)}v n_0,
    \label{eq:spinbrown_BDN_q2_rate2}\\
    (u,v)&\longrightarrow(u-2,v+1),
    &
    \Gamma_{-}^{(\perp)}
    &=
    4\sigma_2^{(2)}u(u-1),
    \label{eq:spinbrown_BDN_q2_rate3}\\
    (u,v)&\longrightarrow(u,v-1),
    &
    \Gamma_{-}^{(z)}
    &=
    8\sigma_2^{(2)}uv.
    \label{eq:spinbrown_BDN_q2_rate4}
\end{align}
The first two transitions increase the ordinary operator size by one and
the last two decrease it by one.

The additional one-site $\sigma^z$ interactions of the full quadratic
model only rotate
\begin{equation}
    \sigma^x\leftrightarrow\sigma^y
\end{equation}
and leave $(u,v)$ unchanged.  Therefore the full and genuine quadratic
models have the same coarse-grained $(u,v)$ master equation even though
their full $k$-twirl generators differ.


\subsection{Large $n$ early-time growth in the two-component model}
\label{subsec:BDN_largeN_growth}

We can now take the same early-time large-$n$ limit used in
\cite{Xu2024}, keeping $u$ and $v$ fixed.  Since
\begin{equation}
    \sigma_q^{(2)}
    =
    \frac{(q-1)!{\cal J}_2}
         {(2n)^{q-1}},
\end{equation}
a transition survives at leading order only if the interaction contains
$q-1$ identity sites and one site already occupied by the operator.

There are two possibilities.

If the occupied site is transverse, the interaction must contain the
opposite transverse Pauli matrix.  The transition is
\begin{equation}
    (u,v)
    \longrightarrow
    (u+q-2,v+1)
\end{equation}
with rate
\begin{equation}
    \Gamma_\perp(u,v)
    =
    4\sigma_q^{(2)}
    2^{q-1}u
    \binom{n-u-v}{q-1}
    \longrightarrow
    4{\cal J}_2u.
    \label{eq:spinbrown_BDN_largeN_perp}
\end{equation}
If the occupied site is a $z$ site, either $x$ or $y$ anticommutes with it.
The transition is
\begin{equation}
    (u,v)
    \longrightarrow
    (u+q,v-1)
\end{equation}
with rate
\begin{equation}
    \Gamma_z(u,v)
    =
    4\sigma_q^{(2)}
    2^qv
    \binom{n-u-v}{q-1}
    \longrightarrow
    8{\cal J}_2v.
    \label{eq:spinbrown_BDN_largeN_z}
\end{equation}
Both processes increase the ordinary size by $q-1$, but their rates are
different.  The early-time dynamics therefore remains intrinsically
two-component even at $n=\infty$.

Because the rates in
Eqs.~\eqref{eq:spinbrown_BDN_largeN_perp} and
\eqref{eq:spinbrown_BDN_largeN_z} are linear in $u$ and $v$, the first
moments obey a closed equation,
\begin{equation}
    \frac{d}{dT}
    \begin{pmatrix}
       \langle u\rangle\\
       \langle v\rangle
    \end{pmatrix}
    =
    4{\cal J}_2
    \begin{pmatrix}
       q-2 & 2q\\
       1   & -2
    \end{pmatrix}
    \begin{pmatrix}
       \langle u\rangle\\
       \langle v\rangle
    \end{pmatrix}.
    \label{eq:spinbrown_BDN_branching_matrix}
\end{equation}
The two eigenvalues of the dimensionless branching matrix are
\begin{equation}
    \lambda_\pm^{(0)}
    =
    \frac{
       q-4
       \pm
       \sqrt{q(q+8)}
    }{2}.
\end{equation}
Thus the dominant early-time exponential rate, with the normalization
Eq.~\eqref{eq:spinbrown_BDN_sigma}, is
\begin{equation}
    \boxed{
    \lambda_{\rm BDN}
    =
    2{\cal J}_2
    \left[
       q-4+\sqrt{q(q+8)}
    \right].
    }
    \label{eq:spinbrown_BDN_lambda}
\end{equation}
In particular,
\begin{equation}
    \lambda_{\rm BDN}
    =
    \begin{cases}
       4{\cal J}_2(\sqrt5-1), & q=2,\\[1mm]
       8\sqrt3\,{\cal J}_2, & q=4.
    \end{cases}
    \label{eq:spinbrown_BDN_lambda_q24}
\end{equation}
Of course an overall Brownian time rescaling changes these numerical
values.  The invariant statement is not the overall coefficient but the
fact that the large-$n$ branching problem is a $2\times2$ problem rather
than the single-rate process of the three-component model.


\subsection{The $\Gamma^3$ charge}
\label{subsec:spinbrown_BDN_charge}

The two-component model also retains the exact charge
\begin{equation}
    \Gamma^3
    =
    \prod_{r=1}^n\sigma_r^z.
\end{equation}
Each constituent operator $O_i$ anticommutes with $\Gamma^3$.  Therefore
\begin{equation}
    \Gamma^3V_I
    =
    (-1)^qV_I\Gamma^3.
\end{equation}
For the quadratic and quartic models,
\begin{equation}
    [\Gamma^3,V_I]=0.
    \label{eq:spinbrown_BDN_charge_commute}
\end{equation}
Hence each replica charge satisfies
\begin{equation}
    [\Gamma^{3(a)},W_I^{(k)}]=0
\end{equation}
and belongs to the commutant of the Brownian Lindblad generators.  The
Brownian semigroup therefore cannot approach the unrestricted Haar twirl
on the full Hilbert space.

This conservation law is also visible directly in
Eq.~\eqref{eq:spinbrown_BDN_master_uv}.  Since
\begin{equation}
    \Gamma^3 S(\Gamma^3)^{-1}
    =
    (-1)^u S,
\end{equation}
the charge sector of an operator is determined by the parity of $u$.  For
even $q$,
\begin{equation}
    \Delta u
    =
    a+b-c-d
    \equiv q
    \equiv0
    \pmod 2,
\end{equation}
so the exact stochastic dynamics preserves
\begin{equation}
    u\pmod2.
\end{equation}


\subsection{Comparison with the Brownian spin model of Xu}

It is useful to summarize the relation between the two models and
\cite{Xu2024}.

For the Swingle-Winer model, the identification is exact.  The Brownian
interaction ensemble consists of all
\begin{equation}
    3^q\binom Nq
\end{equation}
Pauli strings of size $q$, precisely as in the Brownian spin model of
\cite{Xu2024}.  The general semigroup of Sec.~6 therefore reproduces the
four-replica Hamiltonian of that work when $k=2$, and projection onto the
operator-size sector gives exactly
Eq.~\eqref{eq:spinbrown_SW_master}.  After choosing the Brownian clock used
there, the large-$N$ results are
\begin{equation}
    \lambda=q-1,
    \qquad
    x_{\rm eq}=\frac34,
\end{equation}
together with the full distribution
Eqs.~\eqref{eq:spinbrown_Xu_distribution} and
\eqref{eq:spinbrown_Xu_z}.

The two-component model of \cite{BasuDasNandy2026} is different.  Its
interaction ensemble contains only $x$ and $y$ constituent operators.
Consequently, the first closure property used in \cite{Xu2024} still
holds: multiplying a Pauli string by an interaction produces another
Pauli string, and the diagonal four-replica sector is invariant.  The
second property does not hold: Pauli strings of equal support size but
different $z$ content have different transition rates.  Ordinary size
$m$ must therefore be replaced by the pair $(u,v)$.

Schematically,
\begin{equation}
    \begin{array}{c|c|c}
    & \text{Swingle-Winer} & \text{Basu-Das-Nandy}\\ \hline
    \text{local alphabet}
      & \{x,y,z\}
      & \{x,y\}
      \\[1mm]
    \text{$q$-site interactions}
      & \text{all }3^q
      & \text{all }2^q\text{ genuine terms}
      \\[1mm]
    \text{Brownian $k$-twirl}
      & e^{-TG^{(k)}}
      & e^{-TG^{(k)}}
      \\[1mm]
    \text{diagonal Pauli sector}
      & \text{closed}
      & \text{closed}
      \\[1mm]
    \text{ordinary size $m$}
      & \text{closed}
      & \text{not closed}
      \\[1mm]
    \text{minimal coarse variables}
      & m
      & (u,v)
      \\[1mm]
    \text{even-$q$ global charge}
      & \text{none generically}
      & \Gamma^3
    \end{array}
    \label{eq:spinbrown_comparison_table}
\end{equation}

Thus the results of \cite{Xu2024} apply directly to the Brownian
Swingle-Winer model, including the quadratic and quartic cases.  They do
not apply directly to the two-component Spin SYK model.  The natural
analogue in that case is the two-variable master equation
Eq.~\eqref{eq:spinbrown_BDN_master_uv}.  For the full quartic model,
same-site collision terms correct this equation at finite $n$, but are
subleading in the fixed-$q$ large-$n$ limit.  The leading Brownian
large-$n$ dynamics of Spin-SYK$_4$ and gSpin-SYK$_4$ are therefore the
same, while their exact finite-$n$ semigroups are different.


   \section{Semigroup in Brownian SYK}

    The Brownian Sachdev-Ye-Kitaev (SYK) model is a variant of the standard SYK
    model where the static quenched disorder terms in the Hamiltonian are
    replaced by time-dependent stochastic variables \cite{saad2018semiclassical}.
    The Hamiltonian is given by:
    \begin{equation}
        \label{SYKHam}H(t) = i^{q/2}\sum_{i_1<\cdots <i_q}J_{i_1\cdots i_q}(t) (\chi
        _{i_1}\cdots \chi_{i_q})\ .
    \end{equation}

    The couplings $J_{I}(t)$ are independent Gaussian white noise processes defined
    by:
    \begin{align}
        \overline{J_I(t)}            & = 0, \nonumber                      \\
        \overline{J_I(t) J_{I'}(t')} & = \sigma\delta_{II'}\delta(t-t')\nn \\
        \sigma                       & =\frac{2^{q-1}q!}{q^{2}N^{q-1}}J\ .
    \end{align}

    Plugging (\ref{VSYK}) into the equation for $\hat{\mathbb{G}}_{eff}^{(k)}$,
    we obtain Lindblad operators
    \begin{eqnarray}
        &&\mathbb{X}_{i_1<\cdots <i_q}=\sum_a (\psi_{i_1}^{(a)}\cdots \psi_{i_q}^{(a)})=\sum_a
        (\chi_{i_1}^{(a)}\cdots \chi_{i_q}^{(a)})
    \end{eqnarray}
    and
    % The diagonal terms in the sum above, i.e. $a=b$, simplify to
    % \begin{align}
    % &\frac{2}{q!}\sum_{\text{distinct} \: i_1\cdots i_q}\sum_{a=1}^k\lb 1-\psi^{(aa)}_{i_1;LR}\cdots \psi^{(aa)}_{i_q;LR}\rb \nonumber\\
    % &=2\lb {N \choose q}-\sum_{i_1<\cdots<i_q}\sum_{a=1}^k\psi^{(aa)}_{i_1;LR}\cdots \psi^{(aa)}_{i_q;LR}\rb.\label{eq:diagonal_terms}
    % \end{align}
    % To compute the off-diagonal terms, i.e., $a\neq b$, we first note that

    % ----

    % Bilinear operators $\psi^{(ab)}_{ll';\alpha\beta}$ provide an orthonormal basis for the set of all operators in the full Hilbert space of $k N$ qubits:
    % \begin{eqnarray}
    % &&\psi^{(ab)}_{ll';\alpha\beta}\equiv \psi^{(a)}_{l \alpha}\psi^{(b)}_{l' \beta}\ .
    % \end{eqnarray}

    % %-----

    % Note that there is asymmetry between the $L$ and $R$ copies in the dictionary between the complex fermions and the qubits. Since we are ultimately interested in the Weyl fermions,
    % Note the asymmetry between $L$ and $R$.
    % As a result,
    % \begin{eqnarray}
    % &&\chi_{2k-1,L}=(-1)^{k-1}Q_{k-1}X_{kL},\qquad \chi_{2k,L}=(-1)^{k-1}Q_{k-1}Y_{kL}\nn\\
    % &&\chi_{2k-1,R}=(-1)^{k-1}Q_{k-1}(-Z_k)Y_k,\qquad \chi_{2k,R}=(-1)^{k-1}Q_{k-1}(-Z_k)X_k\nn\\
    % &&Z_{kL}=-i\chi_{2k-1,L}\chi_{2k,L},\qquad Z_{kR}=i\chi_{2k-1,R}\chi_{2k,R}\nn\\
    % &&Q_k=\prod_{j=1}^{2k}
    % \end{eqnarray}

    % Written in terms of the Weyl fermions, we find
    % \begin{eqnarray}
    % &&(\chi_{2k-1,L}+i\chi_{2k-1,R})\ket{e}^{\otimes N/2}=0=(\chi_{2k,L}+i\chi_{2k,R})\ket{e}^{\otimes N/2}\nn\\
    % && (\chi_{2k-1,L}-i\chi_{2k-1,R})\ket{f}^{\otimes N/2}=0=(\chi_{2k,L}+i\chi_{2k,R})\ket{f}^{\otimes N/2}\ .
    % \end{eqnarray}
    % Note that $(\chi_{2k,L}+i\chi_{2k-1,R})$ kills both vectors above. Notice that the vector $\ket{e}^{\otimes N/2}$ has the advantage that we do not need to distinguish between $\chi_{2k-1}$ and $\chi_{2k}$, we can simply write the mirror equation
    % \begin{eqnarray}
    % \chi_{j,L}\ket{e}^{\otimes N/2}=-i\chi_{j,R}\ket{e}^{\otimes N/2}\ .
    % \end{eqnarray}

    % In the GNS Hilbert space, the Hamiltonian of SYK for $q$ and $N$ both even can be written as
    % \begin{eqnarray}
    % \hat{H}=H_L-H_R=\sum_{i_1<\cdots<i_q}J_{i_1\cdots i_q}\lb \chi_{i_1L}\cdots \chi_{i_qL}-  \chi_{i_1R}\cdots \chi_{i_qR}\rb
    % \end{eqnarray}
    % Therefore, we can write $\hat{H}=H\otimes 1-1\otimes H$.

    % -----

    % We define the complex fermions
    % \begin{eqnarray}
    % \hat{\chi}_j^{(a)}=\chi^{(a),L}_j-i\chi^{(a),R}_j\ .
    % \end{eqnarray}
    % and we denote their corresponding vacuum with the infinite temperature TFD vector $\ket{1}$ is defined as
    % \begin{eqnarray}
    % \hat{\chi}^{(a)}_j\ket{1}^{(a)}=0\ .
    % \end{eqnarray}

    % Define the operators
    % \begin{eqnarray}
    % Q^{(a)}=\chi^{(a)}_1\cdots \chi^{(a)}_N
    % \end{eqnarray}
    % so that
    % \begin{eqnarray}
    % Q^{(a)}\chi_j^{(a)}=(-1)^{N-j}\chi^{(a)}_1\cdots \chi_{j-1}^{(a)}\chi^{(a)}_j\cdots \chi_N^{(a)}=(-1)^{N-1}\chi_j^{(a)}Q^{(a)}\ .
    % \end{eqnarray}
    % We assume that $N$ is even, so that
    % \begin{eqnarray}
    % &&Q^{(a)}\chi^{(j)}=-\chi_j^{(a)}Q^{(a)},\qquad Q^{(a)}Q^{(a)}=(-1)^{(N/2 \mod 2 )}\ .
    % \end{eqnarray}
    % Following \cite{}, we define
    % \begin{eqnarray}
    % \psi^{(a)}_j=i Q\chi^{(a)}_j
    % \end{eqnarray}
    % so that
    % \begin{eqnarray}
    % \chi^{(a)}_{i_1}\cdots \chi^{(a)}_{i_M}=(-1)^{N M}\psi^{(a)}_{i_1}\cdots \psi^{(a)}_{i_M}
    % \end{eqnarray}
    % where we are assuming both $N$ and $M$ are even.

    % For example, for the case $q=2$, the Lindblad operators are
    % \begin{eqnarray}
    % \mathbb{X}
    % \end{eqnarray}

    % We will see that thanks to the Brownian time-dependent noise, the averaged dynamical evolution of the system above is exactly solvable even at finite $N$. The time-evolution operator is
    % \begin{eqnarray}
    %     &&U(t)=\mathcal{T}\lb e^{-i\int_0^t dt' H(t')}\rb
    % \end{eqnarray}
    % and we are interested in computing
    % \begin{eqnarray}
    %     &&\overline{U(T)\otimes U(T)^\dagger}=\overline{\mathcal{T}\lb e^{-i\int_0^t dt\: \tilde{H}(t')}\rb}\nn\\
    %     &&\tilde{H}(t)=H_L(t)-H_R(t)=i^{q/2}\sum_I J_I(t)(\Psi_I^L-\Psi_I^R),
    % \end{eqnarray}
    % where $H_L$ and $H_R$ are two identical copies of $H$.

    % The Brownian averages over $J(t)$ can be done using a path-integral
    % \begin{eqnarray}
    %     &&\int \mathcal{D}J_I\: e^{-\frac{1}{2\sigma}\sum_I\int dt' J_I(t')^2}\ .
    % \end{eqnarray}
    % Therefore,
    % \begin{eqnarray}
    %     \overline{U(T)\otimes U(T)^\dagger}&&=\int\mathcal{D}J_I \exp\lb\frac{-1}{2\sigma}\sum_I \int_0^T dt'dt J_I(t')\delta(t'-t)\lb J_I(t)+2i\sigma \mathbb{X}_I\rb\rb\nn\\
    %     &&=e^{-\frac{\sigma T}{2}\sum_I \mathbb{X}_I^2}\int \mathcal{D}\tilde{J}_I\exp\lb \frac{-1}{2\sigma}\sum_I \tilde{J}_I^2\rb\nn\\
    %     &&=e^{-\frac{\sigma T}{2}\sum_I \mathbb{X}_I^2}\nn\\
    %     &&\tilde{J}=J+i\sigma \mathbb{X}_I\ .
    % \end{eqnarray}
    % In other words, we find
    % \begin{eqnarray}
    %    &&\overline{U(T)\otimes U(T)^\dagger}=e^{-T H_{eff}} \nn\\
    %    &&H_{eff}=\frac{\sigma}{2}\tilde{H}^2=\frac{\sigma(-1)^q}{2}\sum_I (\Psi^L_I-\Psi^R_I)^2\ .
    % \end{eqnarray}
    % This argument can be repeated for any $H=\sum_I J_I(t) H_I$. Since
    % \begin{eqnarray}
    %     \tilde{H} X\ket{1}=\sum_I [H_I,\mathcal{O}]\ket{1}\nn\ .
    % \end{eqnarray}
    % The super-operator that corresponds to $\tilde{H}_{eff}$ is
    % \begin{eqnarray}
    %     \sum_I [H_I,[H_I,\mO]]\ .
    % \end{eqnarray}

    \subsection{Time-dependent correlators}

    We can compute time-dependent correlators as well. The first observation is that
    we can always order times such that $t_{1}<t_{2}<\cdots <t_{k}$. Then, it follows
    from the calculation in the previous section that
    \begin{eqnarray}
        \mathcal{E}_g\lb\otimes_{i=1}^k \lb U(t_i)\otimes U(t_i)^*\rb\rb=e^{-(t_k-t_{k-1}) \mathbb{G}^{(1)}_{eff}}e^{-(t_2-t_1)\mathbb{G}^{(k-1)}_{eff}}e^{-t_1\mathbb{G}^{(k)}_{eff}}
    \end{eqnarray}
    where $G_{eff}^{(j)}$ is defined using (\ref{eq:Geff_sum_over_AB}) with $\alpha
    ,\beta=L,R$ but $a,b=1,\cdots j$.

    We would like to use this to compute correlation functions

    % \NL{What is the $k$-fold of the semigroup? expand in small $t$: $1-t\mathcal{L}$, and construct the probability measure over unitaries in this limit.}



    \section{Approaching \texorpdfstring{$k$}{k}-designs}

    We are interested in studying how quickly our Brownian dynamics approaches
    $k$-design. In other words, we would like to compare the $k$-th moment of
    our Brownian dynamics
    \begin{eqnarray}
        \mathcal{E}^{(k)}_{g;T}(\mathbb{Y})=e^{-T\mathcal{G}^{(k)}}(\mathbb{Y})
    \end{eqnarray}
    to that of a Haar random ensemble
    \begin{eqnarray}
        \mathcal{T}_H^{(k)}(\mathbb{Y})=\int d\mu_{Haar}(U)\:\lb U^{\otimes k}(\mathbb{Y})(U^\dagger)^{\otimes k}\rb
    \end{eqnarray}
    $d\mu(U)$ denotes the Haar measure, and we have added the superscript $(k)$ to
    keep track of the dimensionality of the space of operator $\mathbb{Y}\in \mathcal{L}
    ((\mathbb{C}^{d})^{\otimes k})$.

    \subsection{Haar Twirl}
    \label{subsec:haar_twirl}

    We start by reviewing the explicit form of the $k$-fold Haar twirl channel
    $\mathcal{T}^{(k)}$. By Schur-Weyl duality, the action of the unitary group
    on the tensor product space commutes with the action of the symmetric group
    $S_{k}$ permuting the tensor factors. As a result, the image of
    $\mathcal{T}^{(k)}$ is the algebra spanned by the permutation partial
    isometries $\{P_{\sigma}\}_{\sigma \in S_k}$:
    \begin{eqnarray}
        P_{\sigma}\ket{i_1\cdots i_k}=\ket{i_{\sigma(1)}\cdots i_{\sigma(k)}}\ .
    \end{eqnarray}
    Assuming $d \geq k$, these operators are linearly independent, allowing us to
    expand the twirled channel as
    \begin{eqnarray}
        \mathcal{T}^{(k)}(\mathbb{Y}) = \sum_{\sigma \in S_k} c_\sigma(\mathbb{Y})
        P_\sigma\ .
    \end{eqnarray}
    To determine the coefficients $c_{\sigma}$, we exploit the fact that
    $\mathcal{T}^{(k)}$ is an orthogonal projection under the Hilbert-Schmidt
    inner product, and
    \begin{eqnarray}
        \Tr(P_\tau^\dagger \mathcal{T}^{(k)}(\mathbb{Y}))&&=\int d\mu_{Haar}(U)\Tr\lb
        (U^\dagger)^{\otimes k} P_\tau^\dagger U^{\otimes k}\: \mathbb{Y}\rb\nn\\
        &&=\int d\mu_{Haar}(U)\Tr\lb P_\tau^\dagger \: \mathbb{Y}\rb=\Tr\lb P_\tau^\dagger
        \mathbb{Y}\rb
    \end{eqnarray}
    where we have used $P_{\tau}^{\dagger}U^{\otimes k}=U^{\otimes k}P_{\tau}^{\dagger}$.
    Substituting the expansion for $\mathcal{T}^{(k)}(X)$ yields a linear system
    involving the Gram matrix of the permutation basis, $G_{\tau, \sigma}= \Tr(P_{\tau}
    ^{\dagger}P_{\sigma}) = d^{\#(\tau^{-1}\sigma)}$, where $\#(\pi)$ counts the
    disjoint cycles in $\pi$. Inverting this relation requires the unitary Weingarten
    function, defined as $Wg(\pi, d) = [G^{-1}]_{\pi}$. This yields the explicit
    closed-form expression for the $k$-fold Haar twirl channel:
    \begin{equation}
        \mathcal{T}^{(k)}(\mathbb{Y}) = \sum_{\sigma, \tau \in S_k}Wg(\sigma\tau^{-1}
        , d) \Tr(P_{\tau}^{\dagger}\mathbb{Y}) P_{\sigma}. \label{eq:general_twirl}
    \end{equation}
    The Weingarten functions are discussed extensively in \cite{collins2006}. As
    a concrete example, for the $k=2$ case relevant to standard randomized
    benchmarking, the symmetric group is $S_{2}= \{e, (12)\}$. The Weingarten
    matrix is the inverse of the Gram matrix $\begin{pmatrix}
        d^{2} & d     \\
        d     & d^{2}
    \end{pmatrix}$, given by
    \begin{equation}
        Wg = \frac{1}{d^{2}- 1}
        \begin{pmatrix}
            1       & -d^{-1} \\
            -d^{-1} & 1
        \end{pmatrix}.
    \end{equation}
    Substituting these values into Eq.~\eqref{eq:general_twirl}, we recover the
    standard projection onto the isotropic (Werner) states:
    \begin{equation}
        \mathcal{T}^{(2)}(\mathbb{Y}) = \frac{\Tr(\mathbb{Y}) - d^{-1}\Tr(\mathbb{Y}S)}{d^{2}-1}
        \mathbb{I}+ \frac{\Tr(\mathbb{Y}S) - d^{-1}\Tr(\mathbb{Y})}{d^{2}-1}S.
    \end{equation}
    where $S$ is the swap operator.

    \subsection{Distance measures}

    To compare our Brownian dynamics to the Haar averages, we use three different
    distance measures:
    \begin{enumerate}
        \item {\bf Hilbert-Schmidt (Frobenius) distance:} This measure is the $(2
            \to 2)$-norm defined of the difference of our superoperators as
            \begin{eqnarray}
                \delta^{HS}_k(T)\equiv \|\mathcal{E}_{g;T}^{(k)}-\mT^{(\otimes)}_H\|_{2\to 2}=\sup_{\mathbb{Y}\neq 0}
                \frac{\|\mathcal{E}_{g;T}^{(k)}(\mathbb{Y})-\mT^{(\otimes)}_{H}(\mathbb{Y})\|_{2}}{\|\mathbb{Y}\|_{2}}
            \end{eqnarray}
            where $\|\mathbb{Y}\|_{2}=\sqrt{\Tr(\mathbb{Y}^{\dagger}\mathbb{Y})}$.
            We will see that for our semigroup, the existence of a gap $\gamma>0$
            in the spectrum of the generator implies that the Frobenius norm decays
            exponentially in time $\gamma T$.

        \item {\bf Diamond distance:} This is the most standard distance measure
            in comparing finite dimensional quantum channels:
            \begin{eqnarray}
                \delta_k^{\diamond}(T)\equiv\sup_m\|(\mathcal{E}_{g;T}^{(k)}-\mT^{(\otimes)}_H)\otimes
                1_m\|_{1\to 1}=\sup_{\mathbb{Y}>0,m}
                \frac{\|((\mathcal{E}_{g;T}^{(k)}-\mT^{(\otimes)}_{H})\otimes 1_{m})(\mathbb{Y})\|_{1}}{\|\mathbb{Y}\|_{1}}
            \end{eqnarray}

        \item {\bf Relative entropy distance:} The relative entropy distance is defined
            as
            \begin{eqnarray}
                \delta^{RE}_k(T)\equiv\sup_m\sup_\rho D((\mathcal{E}^{(k)}_{g;T}\otimes
                1_m)(\rho)\|(\mathcal{T}^{(k)}_H\otimes 1_m)(\rho))
            \end{eqnarray}
            This type of distance measure is often studied using logarithmic Sobolev
            inequalities.

        \item {\bf $\ep$-Relative errors:} The strongest distance measure we consider
            is
            \begin{eqnarray}
                \ep_k(T)\equiv \inf_\ep\{\ep>0:(1-\ep)(\mathcal{E}^{(k)}_{g;T}\circ
                \mathcal{T}^{(k)}_H)\leq \mathcal{E}^{(k)}_H\leq (1+\ep)(\mathcal{E}^{(k)}_{g;T}\circ
                \mathcal{T}^{(k)}_H)\:\}
            \end{eqnarray}
            If all permutations remain invariant under $\mathcal{E}^{(k)}_{g;T}$,
            then our distance measure simplifies to
            \begin{eqnarray}
                \ep_k(T)\equiv \inf_\ep\{\ep>0:(1-\ep)\mathcal{T}^{(k)}_H\leq \mathcal{E}^{(k)}_H\leq
                (1+\ep)\mathcal{T}^{(k)}_H\:\}
            \end{eqnarray}
    \end{enumerate}
    Note that a priori, it is not clear that any of the distance measures above has
    a large $T$ limit in our dynamical system of interest.

    \subsection{Steady states of Brownian dynamics}

    The steady states of our Brownian dynamics are those that are killed by the
    Lindbladian generator $\mathcal{G}_{eff}^{(k)}$.
    \begin{equation}
        \mathcal{G}_{eff}^{(k)}(\mathbb{Y}) = -\sum_{I}\left( \mathbb{X}_{I}\mathbb{Y}
        \mathbb{X}_{I}- \frac{1}{2}\{ \mathbb{X}_{I}^{2}, \mathbb{Y}\} \right) =
        \frac{1}{2}\sum_{I}[\mathbb{X}_{I}, [\mathbb{X}_{I}, \mathbb{Y}]]. \label{eq:generator_form}
    \end{equation}
    For simplicity, we have absorbed the variance $\sigma$ (cf.\ its definition
    in the Brownian ensemble) into the definition of $\mathbb{X}_{I}$. We define
    the \textit{spectral gap} $\Delta_{k}$ as the smallest real part of the non-zero
    eigenvalues of $\mathcal{G}^{(k)}$. The following theorem governs the
    convergence of the dynamics to the Haar average.

    \begin{theorem}[Convergence to steady states]
        Let $\mathfrak{L}\subseteq \mathfrak{su}(d)$ be the Lie algebra generated
        by the set of operators $\{ V_{I}\}_{I}$. The following three statements
        are equivalent:
        \begin{enumerate}
            \item \textbf{Algebraic universality:} The algebra $\mathfrak{L}$
                acts irreducibly on $\mathcal{H}$; i.e.
                $\mathfrak{L}= \mathfrak{u}(d)$.

            \item \textbf{Steady state space:} The kernel of the generator
                coincides exactly with the commutant of the unitary group, i.e.,
                $\ker(\mathcal{G}^{(k)}) = \mathrm{span}\{P_{\sigma}\mid \sigma \in
                S_{k}\}$.

            \item \textbf{Strictly Positive Gap:} The spectral gap satisfies
                $\Delta_{k}> 0$, implying exponential convergence to the Haar twirl
                in any norm:
                \begin{equation}
                    \| e^{-T\mathcal{G}^{(k)}}- \mathcal{T}^{(k)}\|_{2\to2}\leq e
                    ^{-\Delta^{(k)} T}.
                \end{equation}
                Since we are working with finite-dimensional matrix algebras, this
                implies convergence in the diamond distance and relative entropy
                distance, as well:
                \begin{eqnarray}
                    \lim_{T\to \infty}\delta^{\diamond}_k(T)=0=\lim_{T\to \infty}\delta^{RE}_k(T)
                \end{eqnarray}
        \end{enumerate}
    \end{theorem}

    \begin{proof}
        We proceed by establishing the chain of implications $(1) \implies (2) \implies
        (3) \implies (1)$.

        \paragraph{$(1) \implies (2)$:}
        Viewing $\mathcal{L}(\mathcal{H}^{\otimes k})$ as a Hilbert space with
        the Hilbert-Schmidt inner product
        $\langle A, B \rangle = \Tr(A^{\dagger}B)$, we have
        \begin{equation}
            \langle \mathbb{Y}, \mathcal{G}^{(k)}(\mathbb{Y}) \rangle = -\frac{1}{2}
            \sum_{I}\| [\mathbb{X}_{I}, \mathbb{Y}] \|_{2}^{2}.
        \end{equation}
        Since norms are non-negative, $\mathcal{G}^{(k)}(\mathbb{Y}) = 0$ if and
        only if $[\mathbb{X}_{I}, \mathbb{Y}] = 0$ for all $I$. This implies that
        $\mathbb{Y}$ must commute with the entire Lie algebra generated by the global
        operators $\{\mathbb{X}_{I}\}$. Since
        $\mathbb{X}_{I}=\sqrt{\sigma}\sum_{a=1}^{k}V_{I}^{(a)}$, this condition is
        equivalent to commuting with the group $G^{\otimes k}$ generated by exponentiating
        the algebra $\mathfrak{L}$. If $\mathfrak{L}= \mathfrak{su}(d)$, then $G
        = SU(d)$. By Schur-Weyl duality, the operators acting on $(\mathbb{C}^{d}
        )^{\otimes k}$ that commute with $U^{\otimes k}$ for all $U \in SU(d)$ are
        precisely the linear combinations of permutation operators $P_{\sigma}$.
        Thus, $\ker(\mathcal{G}^{(k)}) = \text{span}(S_{k})$.

        \paragraph{$(2) \implies (3)$:}
        The superoperator $\mathcal{G}^{(k)}$ viewed as an operator on the GNS
        Hilbert space is $\hat{\mathbb{G}}^{(k)}$
        \begin{eqnarray}
            &&\mathcal{G}^{(k)}(\mathbb{Y})\ket{1}^{\otimes k}=\hat{\mathbb{G}}^{(k)}(\mathbb{Y}\otimes
            1^{\otimes k})\ket{1}^{\otimes k}\nn\\
            &&\hat{\mathbb{G}}^{(k)}=\sum_I \hat{\mathbb{X}}_I^2
        \end{eqnarray}
        which is manifestly positive semi-definite. Therefore, the spectrum of
        the superoperator $\mathcal{G}^{(k)}$ is comprised of non-negative
        eigenvalues $\lambda$'s. Statement (2) implies that the zero eigenspace ($\lambda
        =0$) corresponds exactly to the image of the Haar twirl $\mathcal{T}^{(k)}$.
        Since the dimension is finite, the spectrum is discrete. Therefore, the smallest
        non-zero eigenvalue, defined as the gap
        $\Delta^{(k)}= \min \{ \lambda_{j}\mid \lambda_{j}> 0 \}$, is strictly positive.
        Standard spectral theory implies that the evolution on the orthogonal
        complement of the kernel decays as
        $\| e^{-T\mathcal{G}^{(k)}}|_{\ker^\perp}\| \le e^{-\Delta^{(k)} T}$.

        \paragraph{$(3) \implies (1)$:}
        We prove by contradiction. Suppose (3) statement holds but (1) does not
        hold; i.e., the generated Lie algebra $\mathfrak{L}$ is a proper subalgebra
        of $\mathfrak{su}(d)$. The dynamical group $G = e^{\mathfrak{L}}$ is
        then a proper subgroup of $SU(d)$. The commutant of a subgroup is strictly
        larger than the commutant of the full group (except in trivial cases, which
        are excluded by the assumption of proper subalgebra). That is,
        $\dim(\ker(\mathcal{G}^{(k)})) = \dim(\text{Comm}(G)) > \dim(\text{span}(
        S_{k}))$. This implies there exists at least one operator $O \in \ker(\mathcal{G}
        ^{(k)})$ such that $O \notin \text{span}(S_{k})$. Consequently, the steady
        state of the evolution includes $O$. The distance to the Haar twirl
        $\mathcal{T}^{(k)}$ (which projects out $O$) will never decay to zero, implying
        the effective gap for convergence to $\mathcal{T}^{(k)}$ is zero.
    \end{proof}
    If the Lie algebra generated by $\mathbb{X}_{I}$ is a proper subalgebra of $\mathfrak{su}
    (d)$, the missing generators generate constants of motion, and should be physically
    interpreted as conserved charges.

    

\appendix


    \section{Correlation functions and $k$-twirl maps}

    Let $\mH$ be the microscopic Hilbert space (dimension $d$) and $\mA=\mathrm{End}
    (\mH)$ the operator algebra. We use the doubled (Choi/GNS) Hilbert space
    \begin{equation}
        \mH_{GNS}:=\mH_{L}\otimes \mH_{R}, \qquad \ket{1}:=\frac{1}{\sqrt{d}}\sum
        _{i=1}^{d}\ket{i}_{L}\ket{i}_{R}.
    \end{equation}
    We identify an operator $A\in\mA$ with the vector
    $\ket{A}:=(A_{L}\otimes 1_{R})\ket{1}$. With this choice, for all
    $A,B\in\mA$,
    \begin{equation}
        \braket{1|A_L\otimes B_R|1}=\frac{1}{d}\tr\big(A\,B^{\mathsf{T}}\big),
    \end{equation}
    where $\mathsf{T}$ denotes transpose in the basis used to define $\ket{1}$.

    \subsection*{Permutation action and cycle states}

    Consider $k$ copies of the GNS Hilbert space, $\mH_{GNS}^{\otimes k}$, with
    canonical vector $\ket{1^{\otimes k}}$. We choose the representation of the permutation
    group $S_{k}$ which acts on the left factors:
    \begin{equation}
        \mathbb{S}_{\pi}\Big(\bigotimes_{a=1}^{k}\ket{i_a}_{L}\ket{j_a}_{R}\Big)
        \,=\,\bigotimes_{a=1}^{k}\ket{i_{\pi(a)}}_{L}\ket{j_a}_{R},\qquad \pi\in
        S_{k}.
    \end{equation}
    For a $k$-cycle $\pi\in C_{k}\subset S_{k}$ and any operators
    $\mO_{1},\dots, \mO_{k}\in\mA$ one has the standard identity
    \begin{equation}
        \bra{\mathbb{S}_\pi 1^{\otimes k}}\Big(\prod_{a=1}^{k}(\mO_{a})_{L}^{(a)}
        \Big)\ket{1^{\otimes k}}\,=\,\tr\big(\mO_{\pi(k)}\cdots \mO_{\pi(1)}\big)
        .
    \end{equation}
    In particular, for the cyclic shift $\pi_{0}=(1\,2\,\dots\,k)$ this gives a cyclic
    trace.

    \subsection*{Complexified evolution as a superoperator}

    We work with dimensionless variables (set $\beta=1$), so that the thermal state
    is $\rho\equiv\rho_{\beta=1}=e^{-H}/Z$ with $Z=\mathrm{Tr}(e^{-H})$.
    Introduce a complexified time
    \begin{equation}
        z=\theta+i t,\qquad \theta\in[0,\tfrac12],\quad t\in\mathbb{R},
    \end{equation}
    and define
    \begin{equation}
        U_{z}:=e^{-zH}.
    \end{equation}
    The corresponding completely-positive ``thermalized'' adjoint action is
    \begin{equation}
        \boxed{\;\mathcal{U}_{z}(\mO):=U_{z}\,\mO\,U_{z}^{\dagger} =e^{-(\theta+i t)H}\,\mO\,e^{-(\theta-i t)H}.\;}
    \end{equation}
    (Note that $\theta\ge 0$ produces thermal damping rather than growth.)

    On the doubled space, left- and right-multiplication are represented by
    \begin{equation}
        \mathbb{U}_{z}:=U_{z,L}\otimes \big(U_{z}^{\dagger}\big)_{R}^{\mathsf{T}}
        =e^{-zH_L}\otimes e^{-z^{\ast} H_R^{\mathsf{T}}}.
    \end{equation}
    For $\vec z=(z_{1},\dots,z_{k})$ define
    \begin{equation}
        \mathbb{U}_{\vec z}:=\bigotimes_{a=1}^{k}\mathbb{U}_{z_a}^{(a)}.
    \end{equation}
    Then the trace of a product of superoperators has the GNS/cycle
    representation
    \begin{equation}
        \boxed{\; \tr\big(\mathcal{U}_{z_1}(\mO_1)\cdots \mathcal{U}_{z_k}(\mO_k)\big) \,=\,\bra{\mathbb{S}_{\pi_0}1^{\otimes k}}\,\mathbb{U}_{\vec z}\,\bigotimes_{a=1}^{k}(\mO_a)_{L}\ket{1^{\otimes k}}.\;}
    \end{equation}
    Thermal correlators arise by imposing the ``thermal strip'' constraint
    \begin{equation}
        \boxed{\;\sum_{a=1}^k \Re z_a=\tfrac12\;}
    \end{equation}
    (which guarantees $\tr(\mathcal{U}_{z_1}(1)\cdots \mathcal{U}_{z_k}(1))=\tr (
    e^{-H})=Z/d$).

    \subsection*{Thermal one- and two-point functions}

    \paragraph{One-point.}
    For any operator $\mO$ and any $t\in\mathbb{R}$,
    \begin{equation}
        \label{1ptthermal}\boxed{\; \tr(\rho\,\mO) =\frac{\bra{1}\,\mathbb{U}_{\frac{1}{2}+i t}\,(\mO)_{L}\ket{1}}{\bra{1}\,\mathbb{U}_{\frac{1}{2}+i t}\ket{1}} =\frac{\tr\big(e^{-\frac{1}{2}H}e^{-i tH}\mO\,e^{+i tH}e^{-\frac{1}{2}H}\big)}{\tr(e^{-H})}=\frac{\tr\big(e^{-H}\mO\big)}{\tr(e^{-H})}\;}
    \end{equation}
    The ratio is $t$-independent by cyclicity of the trace.

    \paragraph{Two-point (Wightman/KMS form).}
    Fix $\tau\in[0,\tfrac12]$ and define the thermal Wightman function
    \begin{equation}
        \label{2ptthermal}G_{12}(t,\tau):=\frac{\tr\big(e^{-(1-\tau)H}\,\mO_{1}(t)\,e^{-\tau
        H}\,\mO_{2}\big)}{\tr(e^{-H})}, \qquad \mO_{1}(t):=e^{+iHt}\mO_{1}e^{-iHt}
        .
    \end{equation}
    It can be written using $k=3$ copies (with the third insertion taken to be the
    identity) as
    \begin{equation}
        \boxed{\; G_{12}(t,\tau) =\frac{\bra{\mathbb{S}_{\pi_0}1^{\otimes 3}}\,\mathbb{U}_{(z_1,z_2,z_3)}\,\big((\mO_{1})_{L}\otimes (\mO_{2})_{L}\otimes 1_{L}\big)\ket{1^{\otimes 3}}}{\bra{\mathbb{S}_{\pi_0}1^{\otimes 3}}\,\mathbb{U}_{(z_1,z_2,z_3)}\ket{1^{\otimes 3}}} ,}
    \end{equation}
    with the convenient choice
    \begin{equation}
        (z_{1},z_{2},z_{3})=(0,\,\tau,\,\tfrac12-\tau).
    \end{equation}
    (Real-time dependence may equivalently be kept in $\mO_{1}(t)$ as above, or
    absorbed into $z_{2}\to \tau+i t$.)

    \subsection*{General $(k\! -\! 1)$-point functions}

    For $k\ge 2$ we define a thermal $(k-1)$-point function by inserting the
    ideneity operator $1$ as the $k$-th operator:
    \begin{equation}
        \boxed{\; G_{1\cdots (k-1)}(\vec z) :=\frac{\bra{\mathbb{S}_{\pi_0}1^{\otimes k}}\,\mathbb{U}_{\vec z}\,\big((\mO_{1})_{L}\otimes\cdots\otimes (\mO_{k-1})_{L}\otimes 1_{L}\big)\ket{1^{\otimes k}}}{\bra{\mathbb{S}_{\pi_0}1^{\otimes k}}\,\mathbb{U}_{\vec z}\ket{1^{\otimes k}}},\qquad \sum_{a=1}^k\Re z_a=\tfrac12.\;}
    \end{equation}
    A concrete parametrization adapted to Euclidean insertion times $\tau_{1},\cdots
    \tau_{k-1}$ with $0=\Re\tau_{1}\le \Re\tau_{2}\le\cdots\le \Re\tau_{k-1}\le \tfrac
    12$ (which fixes the overall time-translation along the thermal circle) is
    \begin{align}
         & G_{1\cdots (k-1)}(\tau_{1},\cdots, \tau_{k-1})\equiv\braket{\mO_1(\tau_1)\cdots \mO_{k-1}(\tau_{k-1})}=G_{1\cdots(k-1)}(\vec{z})\nn \\
         & \Re\tau_{1}=\Re z_{1},\qquad \Re\tau_{k}=1/2-\Re\tau_{k-1}\nn                                                                       \\
         & \forall 1<j<k:\qquad \Re\tau_{j}=\sum_{a=1}^{j}(\Re z_{a}+\Re z_{a+1}),\qquad \tau_{j}=\Re \tau_{j}+i\Im z_{j}
    \end{align}

\section{Conserved charges in Brownian SYK}

    We would like compute the commutant of the Brownian $k$-twirl generators. Let
    \[
        \mathrm{Comm}\{\mathbb{X}_{I}\}\;:=\;\Big\{Y\in\mathrm{End}(\mathcal{H}^{\otimes
        k})\;:\;[\mathbb{X}_{I},Y]=0\ \ \forall I\Big\}, \qquad \mathbb{X}_{I}=\sqrt{\sigma}
        \sum_{a=1}^{k}V_{I}^{(a)}.
    \]
    Since different replicas commute, the Lie algebra generated by
    $\{\mathbb{X}_{I}\}$ is the diagonal action of the single-copy dynamical
    group $G:=\langle e^{itV_I}\rangle$, i.e.
    \[
        \mathrm{Comm}\{\mathbb{X}_{I}\}=\{Y:\ [Y,U^{\otimes k}]=0\ \forall U\in G
        \},
    \]
    the centralizer algebra of the diagonal representation.

    \paragraph{Explicit form for Brownian SYK (even $q$).}
    For SYK with even $q$, each
    $V_{I}^{(a)}=i^{q/2}\chi^{(a)}_{i_1}\cdots\chi^{(a)}_{i_q}$ (equivalently with
    $\psi$ in place of $\chi$) is an even fermion operator and commutes with the
    parity on that replica,
    \[
        Q^{(a)}:=(-i)^{N/2}\prod_{j=1}^{N}\chi^{(a)}_{j}\quad (\text{or }(-i)^{N/2}
        \prod_{j}\psi^{(a)}_{j}),\qquad [Q^{(a)},\mathbb{X}_{I}]=0\ \ \forall a,
        I.
    \]
    Moreover replica permutations $P_{\sigma}$ commute with $\mathbb{X}_{I}$ because
    they relabel the replica sum:
    $P_{\sigma}\mathbb{X}_{I}P_{\sigma}^{-1}=\mathbb{X}_{I}$. Hence always
    \[
        \mathrm{Alg}\langle\,P_{\sigma},\ Q^{(1)},\dots,Q^{(k)}\,\rangle\ \subseteq
        \ \mathrm{Comm}\{\mathbb{X}_{I}\}.
    \]
    If the SYK generators $\{V_{I}\}$ act irreducibly on each fixed-parity sector
    (equivalently, generate the full even operator algebra on one copy), then this
    inclusion is an equality and a convenient linear basis is
    \[
        \boxed{\quad \mathrm{Comm}\{\mathbb{X_I}\} =\mathrm{span}\Big\{\,P_\sigma\,\prod_{a=1}^k (Q^{(a)})^{\epsilon_a}\ :\ \sigma\in S_k,\ \epsilon_a\in\{0,1\}\Big\}. \quad}
    \]
    Equivalently, using parity projectors
    $\Pi_{\vec s}:=\prod_{a=1}^{k}\frac{1+s_{a}Q^{(a)}}{2}$ with $s_{a}=\pm1$, one
    may take the basis $\{\Pi_{\vec s}P_{\sigma}\}$. In the common setting where
    each copy is restricted to a fixed parity sector, $Q^{(a)}=\pm1$ is scalar and
    the commutant reduces to $\mathrm{span}\{P_{\sigma}\}$.

    \bibliographystyle{unsrt}
    \bibliography{ref}
\end{document}